\documentclass[11pt]{report}
% Neuron document style: Charter + Noto Serif CJK KR, light outlined boxes, grayscale.
% Usage: \documentclass[11pt]{report} % Neuron document style: Charter + Noto Serif CJK KR, light outlined boxes, grayscale.
% Usage: \documentclass[11pt]{report} % Neuron document style: Charter + Noto Serif CJK KR, light outlined boxes, grayscale.
% Usage: \documentclass[11pt]{report} \input{neuron-style}  (compile with xelatex)
\usepackage[a4paper,margin=24mm,headheight=14pt]{geometry}
\usepackage{fontspec}
\setmainfont{Charter}[Path=fonts/,Extension=.otf,UprightFont=*-Regular,ItalicFont=*-Italic,
  BoldFont=*-Bold,BoldItalicFont=*-BoldItalic,Ligatures=TeX]
\setsansfont{lmsans10-regular.otf}[BoldFont=lmsans10-bold.otf,ItalicFont=lmsans10-oblique.otf]
\setmonofont{lmmono10-regular.otf}[ItalicFont=lmmono10-italic.otf]
\newfontfamily\krfont{Nanum Myeongjo}[Scale=0.9]
\XeTeXinterchartokenstate=1
\newXeTeXintercharclass\krclass
\count255="AC00 \loop\ifnum\count255<"D7A4 \XeTeXcharclass\count255=\krclass \advance\count255 by 1 \repeat
\count255="3000 \loop\ifnum\count255<"3040 \XeTeXcharclass\count255=\krclass \advance\count255 by 1 \repeat
\count255="3130 \loop\ifnum\count255<"3190 \XeTeXcharclass\count255=\krclass \advance\count255 by 1 \repeat
\count255="1100 \loop\ifnum\count255<"1200 \XeTeXcharclass\count255=\krclass \advance\count255 by 1 \repeat
\XeTeXinterchartoks 0 \krclass = {\begingroup\krfont}
\XeTeXinterchartoks 4095 \krclass = {\begingroup\krfont}
\XeTeXinterchartoks \krclass 0 = {\endgroup}
\XeTeXinterchartoks \krclass 4095 = {\endgroup}
\XeTeXlinebreaklocale "ko"
\XeTeXlinebreakskip=0pt plus 1pt
\usepackage{amsmath,amssymb,amsthm,mathtools,bm}
\usepackage[italic,defaultmathsizes,noasterisk]{mathastext}
\emergencystretch=2.5em

\usepackage{graphicx,booktabs,tabularx,longtable,array,multirow,enumitem,xcolor}
\usepackage{caption}
\captionsetup{font=small,labelfont=bf}
\usepackage{tikz}
\usetikzlibrary{shapes.geometric,arrows.meta,positioning,fit,calc,backgrounds,decorations.pathreplacing}
\usepackage[most]{tcolorbox}
\usepackage{listings}
\usepackage{fancyhdr}
\usepackage[hidelinks]{hyperref}

\setlist{itemsep=1pt,topsep=3pt,parsep=0pt}
\setlength{\parskip}{3pt}
\renewcommand{\arraystretch}{1.12}

\lstdefinestyle{sql}{language=SQL,basicstyle=\ttfamily\small,keywordstyle=\bfseries,
  commentstyle=\itshape\color{black!60},columns=fullflexible,keepspaces=true,
  frame=l,framerule=0.6pt,rulecolor=\color{black!40},xleftmargin=8pt,framexleftmargin=4pt,
  aboveskip=5pt,belowskip=5pt,showstringspaces=false,breaklines=true,
  morekeywords={OVER,PARTITION,RANK,DENSE_RANK,ROW_NUMBER,NTILE,LAG,LEAD,ROWS,RANGE,PRECEDING,FOLLOWING,UNBOUNDED,CURRENT,ROW,NVL,NVL2,REFERENCES,CASCADE,MINUS,INTERSECT,EXISTS,MATERIALIZED,VIEW,REPLACE,RATIO_TO_REPORT,FIRST_VALUE,LAST_VALUE,INTERVAL,NULLS,LAST,FIRST,CHECK,CONSTRAINT,UNIQUE,DEFERRABLE,COMMIT,BEGIN,ROLLBACK,TRANSACTION,INDEX,DEFAULT,ANY,ALL,SOME,ROUND}}
\lstset{style=sql}

\tcbset{
 nb/.style={enhanced,breakable,colback=white,coltitle=black,fonttitle=\bfseries,
   left=3mm,right=3mm,top=1.6mm,bottom=1.6mm,before skip=9pt,after skip=9pt,parbox=false},
 framed/.style={nb,boxrule=0.5pt,colframe=black!70,arc=1.2mm,colbacktitle=white,
   titlerule=0.4pt,titlerule style={black!25},toptitle=1mm,bottomtitle=0.6mm},
 light/.style={framed,colframe=black!35},
 ruled/.style={nb,boxrule=0pt,leftrule=1.1pt,colframe=black!45,sharp corners,left=3.5mm,
   top=0.5mm,bottom=0.5mm,detach title,before upper={\tcbtitle\par\smallskip}},
}
\DeclareTColorBox{keybox}{o}{framed,leftrule=2.2pt,title={\IfValueTF{#1}{#1}{Key idea}}}
\DeclareTColorBox{formulabox}{o}{framed,IfValueTF={#1}{title={#1}}{}}
\DeclareTColorBox{workout}{o}{light,title={Example\IfValueT{#1}{\ \textmd{---}\ #1}}}
\DeclareTColorBox{qa}{m}{light,title={#1}}
\DeclareTColorBox{warnbox}{o}{ruled,leftrule=2.2pt,colframe=black,title={Caution\IfValueT{#1}{\ \textmd{---}\ #1}}}
\DeclareTColorBox{whybox}{o}{ruled,title={\IfValueTF{#1}{#1}{Why it works}}}
\DeclareTColorBox{prof}{o}{ruled,title={Lecture note\IfValueT{#1}{\ \textmd{---}\ #1}}}
\DeclareTColorBox{career}{o}{ruled,title={Career note\IfValueT{#1}{\ \textmd{---}\ #1}}}
\DeclareTColorBox{insight}{o}{ruled,title={In practice\IfValueT{#1}{\ \textmd{---}\ #1}}}

\newcommand{\Step}[1]{\par\smallskip\noindent\textbf{Step #1.}\ }
\newcommand{\Concept}[1]{\hfill{\small[개념: #1]}}
\newcommand{\Use}[1]{\noindent\textit{사용 개념:} #1\par}
\newcommand{\Ans}[1]{\par\smallskip\noindent\fbox{\parbox{\dimexpr\linewidth-2\fboxsep-2\fboxrule\relax}{\textbf{답}\ \ #1}}\par}
\newcommand{\Trap}[1]{\par\smallskip\noindent\textit{흔한 실수:} #1}
\newcommand{\Verify}[1]{\par\smallskip\noindent\textit{검산:} #1}
\newcommand{\Intuit}[1]{\par\smallskip\textbf{직관: #1}}
\newcommand{\pk}[1]{\underline{#1}}
\newcommand{\NULL}{\textsc{null}}
\newcommand{\sqlk}[1]{\texttt{#1}}
\newcommand{\Nul}{{\color{black!50}\textit{null}}}

% relational algebra symbols
\newcommand{\sel}{\sigma}
\newcommand{\proj}{\pi}
\newcommand{\join}{\bowtie}

\pagestyle{fancy}
\fancyhf{}
\fancyhead[R]{\small\itshape\leftmark}
\fancyfoot[C]{\small\thepage}
\renewcommand{\headrulewidth}{0pt}
\renewcommand{\chaptermark}[1]{\markboth{#1}{}}

\usepackage{titlesec}
\titleformat{\chapter}[display]{\bfseries\Large}{\chaptertitlename\ \thechapter}{6pt}{\LARGE}
\titlespacing*{\chapter}{0pt}{-10pt}{14pt}
\titleformat{\section}{\bfseries\large}{\thesection}{0.8em}{}
\titleformat{\subsection}{\bfseries\normalsize}{\thesubsection}{0.8em}{}

% ---------- cheat-sheet / problem-sheet helpers ----------
\newcounter{qn}
\newcommand{\Q}[2]{\par\smallskip\refstepcounter{qn}\noindent\textbf{Q\theqn.}\ {\footnotesize[#1]}\ #2\par}
\newcommand{\A}[1]{\vspace{1pt}\noindent\begin{tcolorbox}[ruled,leftrule=1.6pt,colframe=black!45,before skip=1pt,after skip=4pt,top=0.5mm,bottom=0.5mm]\small\textbf{답} #1\end{tcolorbox}}
\newcommand{\Lab}[3]{\item \textbf{#1}\\ \textit{해 본 것:} #2\\ \textit{\ensuremath{\Rightarrow} 개념:} #3}
\newcommand{\eg}{\textit{예)}\ }
\newcommand{\bi}{\textit{비유)}\ }
  (compile with xelatex)
\usepackage[a4paper,margin=24mm,headheight=14pt]{geometry}
\usepackage{fontspec}
\setmainfont{Charter}[Path=fonts/,Extension=.otf,UprightFont=*-Regular,ItalicFont=*-Italic,
  BoldFont=*-Bold,BoldItalicFont=*-BoldItalic,Ligatures=TeX]
\setsansfont{lmsans10-regular.otf}[BoldFont=lmsans10-bold.otf,ItalicFont=lmsans10-oblique.otf]
\setmonofont{lmmono10-regular.otf}[ItalicFont=lmmono10-italic.otf]
\newfontfamily\krfont{Nanum Myeongjo}[Scale=0.9]
\XeTeXinterchartokenstate=1
\newXeTeXintercharclass\krclass
\count255="AC00 \loop\ifnum\count255<"D7A4 \XeTeXcharclass\count255=\krclass \advance\count255 by 1 \repeat
\count255="3000 \loop\ifnum\count255<"3040 \XeTeXcharclass\count255=\krclass \advance\count255 by 1 \repeat
\count255="3130 \loop\ifnum\count255<"3190 \XeTeXcharclass\count255=\krclass \advance\count255 by 1 \repeat
\count255="1100 \loop\ifnum\count255<"1200 \XeTeXcharclass\count255=\krclass \advance\count255 by 1 \repeat
\XeTeXinterchartoks 0 \krclass = {\begingroup\krfont}
\XeTeXinterchartoks 4095 \krclass = {\begingroup\krfont}
\XeTeXinterchartoks \krclass 0 = {\endgroup}
\XeTeXinterchartoks \krclass 4095 = {\endgroup}
\XeTeXlinebreaklocale "ko"
\XeTeXlinebreakskip=0pt plus 1pt
\usepackage{amsmath,amssymb,amsthm,mathtools,bm}
\usepackage[italic,defaultmathsizes,noasterisk]{mathastext}
\emergencystretch=2.5em

\usepackage{graphicx,booktabs,tabularx,longtable,array,multirow,enumitem,xcolor}
\usepackage{caption}
\captionsetup{font=small,labelfont=bf}
\usepackage{tikz}
\usetikzlibrary{shapes.geometric,arrows.meta,positioning,fit,calc,backgrounds,decorations.pathreplacing}
\usepackage[most]{tcolorbox}
\usepackage{listings}
\usepackage{fancyhdr}
\usepackage[hidelinks]{hyperref}

\setlist{itemsep=1pt,topsep=3pt,parsep=0pt}
\setlength{\parskip}{3pt}
\renewcommand{\arraystretch}{1.12}

\lstdefinestyle{sql}{language=SQL,basicstyle=\ttfamily\small,keywordstyle=\bfseries,
  commentstyle=\itshape\color{black!60},columns=fullflexible,keepspaces=true,
  frame=l,framerule=0.6pt,rulecolor=\color{black!40},xleftmargin=8pt,framexleftmargin=4pt,
  aboveskip=5pt,belowskip=5pt,showstringspaces=false,breaklines=true,
  morekeywords={OVER,PARTITION,RANK,DENSE_RANK,ROW_NUMBER,NTILE,LAG,LEAD,ROWS,RANGE,PRECEDING,FOLLOWING,UNBOUNDED,CURRENT,ROW,NVL,NVL2,REFERENCES,CASCADE,MINUS,INTERSECT,EXISTS,MATERIALIZED,VIEW,REPLACE,RATIO_TO_REPORT,FIRST_VALUE,LAST_VALUE,INTERVAL,NULLS,LAST,FIRST,CHECK,CONSTRAINT,UNIQUE,DEFERRABLE,COMMIT,BEGIN,ROLLBACK,TRANSACTION,INDEX,DEFAULT,ANY,ALL,SOME,ROUND}}
\lstset{style=sql}

\tcbset{
 nb/.style={enhanced,breakable,colback=white,coltitle=black,fonttitle=\bfseries,
   left=3mm,right=3mm,top=1.6mm,bottom=1.6mm,before skip=9pt,after skip=9pt,parbox=false},
 framed/.style={nb,boxrule=0.5pt,colframe=black!70,arc=1.2mm,colbacktitle=white,
   titlerule=0.4pt,titlerule style={black!25},toptitle=1mm,bottomtitle=0.6mm},
 light/.style={framed,colframe=black!35},
 ruled/.style={nb,boxrule=0pt,leftrule=1.1pt,colframe=black!45,sharp corners,left=3.5mm,
   top=0.5mm,bottom=0.5mm,detach title,before upper={\tcbtitle\par\smallskip}},
}
\DeclareTColorBox{keybox}{o}{framed,leftrule=2.2pt,title={\IfValueTF{#1}{#1}{Key idea}}}
\DeclareTColorBox{formulabox}{o}{framed,IfValueTF={#1}{title={#1}}{}}
\DeclareTColorBox{workout}{o}{light,title={Example\IfValueT{#1}{\ \textmd{---}\ #1}}}
\DeclareTColorBox{qa}{m}{light,title={#1}}
\DeclareTColorBox{warnbox}{o}{ruled,leftrule=2.2pt,colframe=black,title={Caution\IfValueT{#1}{\ \textmd{---}\ #1}}}
\DeclareTColorBox{whybox}{o}{ruled,title={\IfValueTF{#1}{#1}{Why it works}}}
\DeclareTColorBox{prof}{o}{ruled,title={Lecture note\IfValueT{#1}{\ \textmd{---}\ #1}}}
\DeclareTColorBox{career}{o}{ruled,title={Career note\IfValueT{#1}{\ \textmd{---}\ #1}}}
\DeclareTColorBox{insight}{o}{ruled,title={In practice\IfValueT{#1}{\ \textmd{---}\ #1}}}

\newcommand{\Step}[1]{\par\smallskip\noindent\textbf{Step #1.}\ }
\newcommand{\Concept}[1]{\hfill{\small[개념: #1]}}
\newcommand{\Use}[1]{\noindent\textit{사용 개념:} #1\par}
\newcommand{\Ans}[1]{\par\smallskip\noindent\fbox{\parbox{\dimexpr\linewidth-2\fboxsep-2\fboxrule\relax}{\textbf{답}\ \ #1}}\par}
\newcommand{\Trap}[1]{\par\smallskip\noindent\textit{흔한 실수:} #1}
\newcommand{\Verify}[1]{\par\smallskip\noindent\textit{검산:} #1}
\newcommand{\Intuit}[1]{\par\smallskip\textbf{직관: #1}}
\newcommand{\pk}[1]{\underline{#1}}
\newcommand{\NULL}{\textsc{null}}
\newcommand{\sqlk}[1]{\texttt{#1}}
\newcommand{\Nul}{{\color{black!50}\textit{null}}}

% relational algebra symbols
\newcommand{\sel}{\sigma}
\newcommand{\proj}{\pi}
\newcommand{\join}{\bowtie}

\pagestyle{fancy}
\fancyhf{}
\fancyhead[R]{\small\itshape\leftmark}
\fancyfoot[C]{\small\thepage}
\renewcommand{\headrulewidth}{0pt}
\renewcommand{\chaptermark}[1]{\markboth{#1}{}}

\usepackage{titlesec}
\titleformat{\chapter}[display]{\bfseries\Large}{\chaptertitlename\ \thechapter}{6pt}{\LARGE}
\titlespacing*{\chapter}{0pt}{-10pt}{14pt}
\titleformat{\section}{\bfseries\large}{\thesection}{0.8em}{}
\titleformat{\subsection}{\bfseries\normalsize}{\thesubsection}{0.8em}{}

% ---------- cheat-sheet / problem-sheet helpers ----------
\newcounter{qn}
\newcommand{\Q}[2]{\par\smallskip\refstepcounter{qn}\noindent\textbf{Q\theqn.}\ {\footnotesize[#1]}\ #2\par}
\newcommand{\A}[1]{\vspace{1pt}\noindent\begin{tcolorbox}[ruled,leftrule=1.6pt,colframe=black!45,before skip=1pt,after skip=4pt,top=0.5mm,bottom=0.5mm]\small\textbf{답} #1\end{tcolorbox}}
\newcommand{\Lab}[3]{\item \textbf{#1}\\ \textit{해 본 것:} #2\\ \textit{\ensuremath{\Rightarrow} 개념:} #3}
\newcommand{\eg}{\textit{예)}\ }
\newcommand{\bi}{\textit{비유)}\ }
  (compile with xelatex)
\usepackage[a4paper,margin=24mm,headheight=14pt]{geometry}
\usepackage{fontspec}
\setmainfont{Charter}[Path=fonts/,Extension=.otf,UprightFont=*-Regular,ItalicFont=*-Italic,
  BoldFont=*-Bold,BoldItalicFont=*-BoldItalic,Ligatures=TeX]
\setsansfont{lmsans10-regular.otf}[BoldFont=lmsans10-bold.otf,ItalicFont=lmsans10-oblique.otf]
\setmonofont{lmmono10-regular.otf}[ItalicFont=lmmono10-italic.otf]
\newfontfamily\krfont{Nanum Myeongjo}[Scale=0.9]
\XeTeXinterchartokenstate=1
\newXeTeXintercharclass\krclass
\count255="AC00 \loop\ifnum\count255<"D7A4 \XeTeXcharclass\count255=\krclass \advance\count255 by 1 \repeat
\count255="3000 \loop\ifnum\count255<"3040 \XeTeXcharclass\count255=\krclass \advance\count255 by 1 \repeat
\count255="3130 \loop\ifnum\count255<"3190 \XeTeXcharclass\count255=\krclass \advance\count255 by 1 \repeat
\count255="1100 \loop\ifnum\count255<"1200 \XeTeXcharclass\count255=\krclass \advance\count255 by 1 \repeat
\XeTeXinterchartoks 0 \krclass = {\begingroup\krfont}
\XeTeXinterchartoks 4095 \krclass = {\begingroup\krfont}
\XeTeXinterchartoks \krclass 0 = {\endgroup}
\XeTeXinterchartoks \krclass 4095 = {\endgroup}
\XeTeXlinebreaklocale "ko"
\XeTeXlinebreakskip=0pt plus 1pt
\usepackage{amsmath,amssymb,amsthm,mathtools,bm}
\usepackage[italic,defaultmathsizes,noasterisk]{mathastext}
\emergencystretch=2.5em

\usepackage{graphicx,booktabs,tabularx,longtable,array,multirow,enumitem,xcolor}
\usepackage{caption}
\captionsetup{font=small,labelfont=bf}
\usepackage{tikz}
\usetikzlibrary{shapes.geometric,arrows.meta,positioning,fit,calc,backgrounds,decorations.pathreplacing}
\usepackage[most]{tcolorbox}
\usepackage{listings}
\usepackage{fancyhdr}
\usepackage[hidelinks]{hyperref}

\setlist{itemsep=1pt,topsep=3pt,parsep=0pt}
\setlength{\parskip}{3pt}
\renewcommand{\arraystretch}{1.12}

\lstdefinestyle{sql}{language=SQL,basicstyle=\ttfamily\small,keywordstyle=\bfseries,
  commentstyle=\itshape\color{black!60},columns=fullflexible,keepspaces=true,
  frame=l,framerule=0.6pt,rulecolor=\color{black!40},xleftmargin=8pt,framexleftmargin=4pt,
  aboveskip=5pt,belowskip=5pt,showstringspaces=false,breaklines=true,
  morekeywords={OVER,PARTITION,RANK,DENSE_RANK,ROW_NUMBER,NTILE,LAG,LEAD,ROWS,RANGE,PRECEDING,FOLLOWING,UNBOUNDED,CURRENT,ROW,NVL,NVL2,REFERENCES,CASCADE,MINUS,INTERSECT,EXISTS,MATERIALIZED,VIEW,REPLACE,RATIO_TO_REPORT,FIRST_VALUE,LAST_VALUE,INTERVAL,NULLS,LAST,FIRST,CHECK,CONSTRAINT,UNIQUE,DEFERRABLE,COMMIT,BEGIN,ROLLBACK,TRANSACTION,INDEX,DEFAULT,ANY,ALL,SOME,ROUND}}
\lstset{style=sql}

\tcbset{
 nb/.style={enhanced,breakable,colback=white,coltitle=black,fonttitle=\bfseries,
   left=3mm,right=3mm,top=1.6mm,bottom=1.6mm,before skip=9pt,after skip=9pt,parbox=false},
 framed/.style={nb,boxrule=0.5pt,colframe=black!70,arc=1.2mm,colbacktitle=white,
   titlerule=0.4pt,titlerule style={black!25},toptitle=1mm,bottomtitle=0.6mm},
 light/.style={framed,colframe=black!35},
 ruled/.style={nb,boxrule=0pt,leftrule=1.1pt,colframe=black!45,sharp corners,left=3.5mm,
   top=0.5mm,bottom=0.5mm,detach title,before upper={\tcbtitle\par\smallskip}},
}
\DeclareTColorBox{keybox}{o}{framed,leftrule=2.2pt,title={\IfValueTF{#1}{#1}{Key idea}}}
\DeclareTColorBox{formulabox}{o}{framed,IfValueTF={#1}{title={#1}}{}}
\DeclareTColorBox{workout}{o}{light,title={Example\IfValueT{#1}{\ \textmd{---}\ #1}}}
\DeclareTColorBox{qa}{m}{light,title={#1}}
\DeclareTColorBox{warnbox}{o}{ruled,leftrule=2.2pt,colframe=black,title={Caution\IfValueT{#1}{\ \textmd{---}\ #1}}}
\DeclareTColorBox{whybox}{o}{ruled,title={\IfValueTF{#1}{#1}{Why it works}}}
\DeclareTColorBox{prof}{o}{ruled,title={Lecture note\IfValueT{#1}{\ \textmd{---}\ #1}}}
\DeclareTColorBox{career}{o}{ruled,title={Career note\IfValueT{#1}{\ \textmd{---}\ #1}}}
\DeclareTColorBox{insight}{o}{ruled,title={In practice\IfValueT{#1}{\ \textmd{---}\ #1}}}

\newcommand{\Step}[1]{\par\smallskip\noindent\textbf{Step #1.}\ }
\newcommand{\Concept}[1]{\hfill{\small[개념: #1]}}
\newcommand{\Use}[1]{\noindent\textit{사용 개념:} #1\par}
\newcommand{\Ans}[1]{\par\smallskip\noindent\fbox{\parbox{\dimexpr\linewidth-2\fboxsep-2\fboxrule\relax}{\textbf{답}\ \ #1}}\par}
\newcommand{\Trap}[1]{\par\smallskip\noindent\textit{흔한 실수:} #1}
\newcommand{\Verify}[1]{\par\smallskip\noindent\textit{검산:} #1}
\newcommand{\Intuit}[1]{\par\smallskip\textbf{직관: #1}}
\newcommand{\pk}[1]{\underline{#1}}
\newcommand{\NULL}{\textsc{null}}
\newcommand{\sqlk}[1]{\texttt{#1}}
\newcommand{\Nul}{{\color{black!50}\textit{null}}}

% relational algebra symbols
\newcommand{\sel}{\sigma}
\newcommand{\proj}{\pi}
\newcommand{\join}{\bowtie}

\pagestyle{fancy}
\fancyhf{}
\fancyhead[R]{\small\itshape\leftmark}
\fancyfoot[C]{\small\thepage}
\renewcommand{\headrulewidth}{0pt}
\renewcommand{\chaptermark}[1]{\markboth{#1}{}}

\usepackage{titlesec}
\titleformat{\chapter}[display]{\bfseries\Large}{\chaptertitlename\ \thechapter}{6pt}{\LARGE}
\titlespacing*{\chapter}{0pt}{-10pt}{14pt}
\titleformat{\section}{\bfseries\large}{\thesection}{0.8em}{}
\titleformat{\subsection}{\bfseries\normalsize}{\thesubsection}{0.8em}{}

% ---------- cheat-sheet / problem-sheet helpers ----------
\newcounter{qn}
\newcommand{\Q}[2]{\par\smallskip\refstepcounter{qn}\noindent\textbf{Q\theqn.}\ {\footnotesize[#1]}\ #2\par}
\newcommand{\A}[1]{\vspace{1pt}\noindent\begin{tcolorbox}[ruled,leftrule=1.6pt,colframe=black!45,before skip=1pt,after skip=4pt,top=0.5mm,bottom=0.5mm]\small\textbf{답} #1\end{tcolorbox}}
\newcommand{\Lab}[3]{\item \textbf{#1}\\ \textit{해 본 것:} #2\\ \textit{\ensuremath{\Rightarrow} 개념:} #3}
\newcommand{\eg}{\textit{예)}\ }
\newcommand{\bi}{\textit{비유)}\ }

\usepackage{float}
\usetikzlibrary{arrows.meta,positioning,calc,decorations.pathreplacing}
\usepackage{pgfplots}\pgfplotsset{compat=1.18}
\newcommand{\ft}[1]{\texttt{\detokenize{#1}}}
\newcolumntype{L}{>{\raggedright\arraybackslash}X}
\lstdefinestyle{py}{language=Python,basicstyle=\ttfamily\small,keywordstyle=\bfseries,commentstyle=\itshape\color{black!55},
  stringstyle=\color{black!80},frame=leftline,framerule=0.8pt,rulecolor=\color{black!45},xleftmargin=3mm,
  showstringspaces=false,breaklines=true,columns=fullflexible,upquote=true}
\lstset{style=py}
\begin{document}

\begin{titlepage}\centering
\vspace*{30mm}
{\Huge\bfseries LendingClub 대출 승인 전략\par}\medskip
{\Large 데이터에서 검증까지, 기초부터 따라가는 강의노트\par}
\vspace{12mm}
{\large 서울대 ABS 14기 팀 프로젝트\par}
\vspace{4mm}
{\small 2026년 11월 9일 기준 \ensuremath{\cdot} 숫자는 모두 공식 train 데이터 실험 결과(예시라고 적은 것 제외)\par}
\vfill
{\small 이 노트는 코드 저장소의 1\textasciitilde 26단계 실험, \ft{docs/decisions.md}, \ft{docs/research_log.md}를 바탕으로 쓴다.\par}
\end{titlepage}

\tableofcontents
\chapter*{How to Read}
\addcontentsline{toc}{chapter}{How to Read}
\markboth{How to Read}{}

이 노트는 세 독자를 함께 생각하고 쓴다.
\begin{itemize}[leftmargin=*]
\item \textbf{나(공부)}: 코드는 돌렸지만 개념이 흐릿한 부분을 기초부터 다시 잡는다. 각 장은 ``질문 \ensuremath{\to} 작은 예시 \ensuremath{\to} 식 \ensuremath{\to} 숫자 대입 \ensuremath{\to} 우리 데이터 결과'' 순서로 간다.
\item \textbf{취업 준비}: 각 장 끝의 \textit{Career note}와 부록 A의 면접 Q\&A는 ``이 프로젝트에서 무엇을 했고 왜 그렇게 했나''를 2분 안에 말할 수 있게 정리한 것이다.
\item \textbf{팀원과 교수님}: 9장은 전체 이야기를 발표 순서로 다시 묶는다. 결론만 필요하면 Lecture Highlights와 9장만 읽어도 된다.
\end{itemize}

\textbf{처음 읽을 때.} 각 핵심 박스 끝의 \textbf{직관:} 줄만 순서대로 읽으면 전체 흐름이 한 번에 잡힌다. 그다음 Example 박스를 손으로 따라 계산한다. 예시 숫자는 모두 \ft{docs/notes/lecture/code/checks.py}로 검산했다.

\textbf{표기.} PD는 부도확률(probability of default), LGD는 부도 시 손실률(loss given default), $r$은 약정 금리, $tb$는 같은 달 같은 만기의 미 국채 수익률이다. ``선별 가치''는 무작위 승인 대비 초과 IRR(\%p)이다. 코드 경로는 \ft{src/lc/...}, 실험은 \ft{scripts/NN_...py}로 적는다.

\chapter*{Glossary}
\addcontentsline{toc}{chapter}{Glossary}
\markboth{Glossary}{}
{\small
\begin{longtable}{@{}p{43mm}p{40mm}p{70mm}@{}}\toprule
English & 한국어 & 한 줄 뜻\\\midrule\endhead
Amortization & 원리금 균등 상환 & 매달 같은 금액으로 이자와 원금을 함께 갚는 방식\\
As-of (point-in-time) training & 시점 학습 & 그날 알 수 있던 정보로만 학습\\
AUC & ROC 곡선 아래 면적 & 부도 하나와 정상 하나를 뽑았을 때 부도에 더 높은 점수를 줄 확률\\
Availability rule & 가용성 규칙 & 학습 시점에 50\% 넘게 비어 있던 변수는 뺀다\\
Block bootstrap & 묶음 부트스트랩 & 서로 얽힌 관측을 묶음(발행월) 단위로 다시 뽑는 방법\\
Brier score & 브라이어 점수 & 예측 확률과 0/1 결과 차이의 제곱 평균, 확률 수준의 정확도\\
Brinson attribution & 브린슨 성과 귀속 & 초과수익을 배분, 선별, 상호작용으로 나눔\\
Calibration & 확률 보정 & 예측 확률을 실제 빈도에 맞추는 일\\
Censoring & 중도 절단 & 관측이 끝날 때 아직 결과가 안 난 상태\\
Charge-off lag & 상각 지연 & 마지막 납부 후 부도 확정까지 걸리는 시간(약 5개월)\\
Cohort & 코호트 & 같은 달에 발행된 대출 묶음\\
Data leakage & 데이터 누수 & 결정 시점에 알 수 없던 정보가 학습에 섞이는 것\\
Data snooping & 데이터 엿보기 & 같은 데이터를 반복해 보며 고르다가 우연에 맞추는 것\\
Early stopping & 조기 종료 & 검증 손실이 더 안 줄면 학습을 멈춤\\
Elastic Net & 엘라스틱 넷 & L1과 L2 벌점을 섞은 회귀\\
Excess return & 초과수익 & 수익률 $-$ 비교 기준 수익률\\
Expected return score & 기대수익 점수 & $(1-p)r - p\cdot\mathrm{LGD} - tb$\\
Gradient boosting & 그래디언트 부스팅 & 앞 나무의 오차를 다음 나무가 고치며 쌓는 앙상블\\
Holdout & 보류 구간 & 마지막 확인용으로 아껴 둔 데이터\\
Holm correction & 홈 보정 & 여러 검정을 할 때 p값을 단계적으로 엄격하게 보정\\
Information ratio & 정보비율 & 기준 대비 초과수익 평균 / 그 표준편차\\
IRR & 내부수익률 & 현금흐름의 현재가치를 0으로 만드는 할인율\\
Information value (IV) & 정보 가치 & 한 변수가 부도/정상을 가르는 힘(WOE의 가중합)\\
Lasso (L1) & 라소 & 계수 절댓값 합에 벌점, 일부 계수를 정확히 0으로 만듦\\
LGD & 부도 시 손실률 & 부도가 나면 원금 중 잃는 비율\\
MAR / MCAR / MNAR & 결측 메커니즘 & 결측 이유가 관측값/무작위/결측된 값 자체에 달려 있음\\
PD & 부도확률 & 대출이 부도날 확률\\
Placebo test & 위약 검정 & 효과가 없어야 할 조건(라벨 섞기)에서 효과가 0인지 확인\\
Pre-registration & 사전 등록 & 결과를 보기 전에 가설과 판정 기준을 고정\\
RAROC & 위험조정 자본수익률 & (수익 $-$ 비용 $-$ 예상손실)/경제적 자본\\
Ridge (L2) & 릿지 & 계수 제곱합에 벌점, 계수를 고르게 줄임\\
Resolution bias & 종결 편향 & 이미 끝난 대출만 보면 빨리 끝난 대출(부도, 조기상환)이 과대 대표됨\\
Selection curve & 선별 곡선 & 승인비율별 (선택 포트폴리오 수익 $-$ 무작위 수익)\\
Selection value & 선별 가치 & 선별 곡선의 평균 높이\\
Sharpe ratio & 샤프 지수 & 초과수익 평균 / 표준편차\\
Survivorship (resolved-only) bias & 종결 표본 편향 & 결과가 난 대출만 남겨서 생기는 왜곡\\
Target encoding & 타깃 인코딩 & 범주를 그 범주의 부도율로 바꿈(누수 주의)\\
Walk-forward & 순차 전진 검증 & 시간 순서대로 학습 기간을 늘리며 다음 기간을 평가\\
WOE & 증거 가중치 & 구간별 $\ln(\text{정상 비중}/\text{부도 비중})$\\\bottomrule
\end{longtable}}

\chapter*{Key Formulas}
\addcontentsline{toc}{chapter}{Key Formulas}
\markboth{Key Formulas}{}
{\small
\begin{longtable}{@{}p{34mm}p{58mm}p{60mm}@{}}\toprule
개념 & 식 & 뜻\\\midrule\endhead
월 상환액 & $A = P\,\dfrac{i}{1-(1+i)^{-T}},\ i=r/12$ & 원금 $P$를 $T$개월에 같은 금액으로 갚는 금액\\[2mm]
IRR & $\sum_{t=0}^{H} \dfrac{CF_t}{(1+m)^t}=0,\ \ \mathrm{IRR}=(1+m)^{12}-1$ & 돈을 넣고 돌려받은 흐름이 실제로 낸 연 수익률\\[2mm]
단순 연환산 & $(\text{받은 총액}/P)^{12/T}-1$ & 재투자를 무시하는 근사(우승 논문식)\\
기대수익 점수 & $s=(1-p)\,r-p\,\mathrm{LGD}-tb$ & 국채보다 얼마나 더 벌 것으로 기대되는가\\
LGD 추정 & $1-\overline{(\text{받은 총액}/P)}_{\text{부도}}$ & 그 시점까지 확정된 부도 대출의 평균 손실\\
Sharpe & $\mathrm{SR}=\bar{x}/\mathrm{sd}(x)$ & 위험 한 단위당 초과수익\\
$\sqrt{a}$ 항등식 & $\mathrm{SR}=\sqrt{a}\,m/\sqrt{s^2+(1-a)m^2}$ & 거절을 0으로 두면 승인비율 $a$에 따라 기계적으로 변함\\
Logistic & $p=\sigma(z)=1/(1+e^{-z}),\ z=\beta_0+\beta^\top x$ & 선형 점수를 확률로\\
L1 / L2 / EN 벌점 & $\lambda\|\beta\|_1,\ \lambda\|\beta\|_2^2,\ \lambda(\alpha\|\beta\|_1+(1-\alpha)\|\beta\|_2^2)$ & 계수를 줄여 과적합 방지\\
WOE, IV & $\mathrm{WOE}_k=\ln\dfrac{g_k/G}{b_k/B},\ \mathrm{IV}=\sum_k(\tfrac{g_k}{G}-\tfrac{b_k}{B})\mathrm{WOE}_k$ & 구간별 증거와 변수 전체의 힘\\[2mm]
Brier & $\frac1n\sum(p_i-y_i)^2$ & 확률 수준의 정확도\\
AUC & $P(p_{\text{부도}}>p_{\text{정상}})$ & 순위의 정확도\\
Brinson & 배분 $\sum(w^p_k-w^b_k)R^b_k$, 선별 $\sum w^b_k(R^p_k-R^b_k)$, 상호작용 $\sum(w^p_k-w^b_k)(R^p_k-R^b_k)$ & 초과수익의 출처 분해\\
선별 가치 & $\frac{1}{9}\sum_{q}\big[\mathrm{IRR}(\text{상위 }q)-\mathrm{IRR}(\text{무작위 }q)\big]$ & 순위가 돈을 얼마나 더 버는가\\
Holm & $p_{(j)}^{adj}=\max_{i\le j}\min\{1,(k-i+1)p_{(i)}\}$ & 여러 비교의 오류율 관리\\\bottomrule
\end{longtable}}

\chapter*{Index of Worked Examples}
\addcontentsline{toc}{chapter}{Index of Worked Examples}
\markboth{Index of Worked Examples}{}
{\small
\begin{longtable}{@{}p{12mm}p{80mm}p{60mm}@{}}\toprule
절 & 예제 & 사용 개념\\\midrule\endhead
\S2.1 & 1만 달러, 연 12\%, 36개월 대출의 월 상환액 & 원리금 균등 상환\\
\S2.2 & 정상 상환 대출의 IRR이 12\%가 아니라 12.68\%인 이유 & IRR, 월 복리\\
\S2.2 & 10개월 내고 부도 난 대출의 IRR과 단순 연환산 & 현금흐름 복원, 손실률\\
\S2.3 & 정상 9건 + 부도 1건 포트폴리오의 IRR & pooled IRR, 수수료\\
\S3.3 & 4건짜리 종결 편향 & resolution bias\\
\S4.2 & 결측 구간의 WOE와 IV & WOE, IV, MNAR\\
\S5.1 & 로지스틱 함수에 숫자 넣기 & sigmoid\\
\S5.6 & 두 대출의 기대수익 점수 비교 & 기대수익, LGD\\
\S6.1 & 승인비율에 따른 Sharpe 변화 & $\sqrt{a}$ 항등식\\
\S6.3 & 5건짜리 AUC와 Brier & 순위 vs 수준\\
\S6.5 & 두 등급 Brinson 분해 & 배분, 선별, 상호작용\\
\S7.3 & 월별 차이 5개의 부트스트랩 구간 & bootstrap\\
\S7.4 & p값 세 개의 Holm 보정 & multiple testing\\\bottomrule
\end{longtable}}

\chapter*{Lecture Highlights}
\addcontentsline{toc}{chapter}{Lecture Highlights}
\markboth{Lecture Highlights}{}

\begin{keybox}[이 프로젝트의 결론 다섯 줄]
\begin{enumerate}[leftmargin=*]
\item \textbf{수익률을 어떻게 재느냐가 답을 바꾼다.} 같은 2013--16년 36개월 대출이 단순 연환산으로는 국채보다 못하고, 현금흐름 IRR로는 48개 발행월 모두 국채보다 높다(IRR 5.73\% vs 0.87\%).
\item \textbf{시간을 지키면 성능이 내려간다.} 무작위 분할 AUC 약 0.73이 시점 학습에서는 0.62--0.68이다. 진짜 성능은 뒤쪽이다.
\item \textbf{기대수익 순위는 무작위보다 약 1.03\%p 더 번다}(95\% 구간 0.95--1.11, 발행월 부트스트랩). LightGBM과 XGBoost가 같은 값을 낸다.
\item \textbf{그 추가 수익은 같은 등급 안에서 나온다.} LC 등급끼리의 가격은 대체로 공정했고, 등급이 덜 반영한 정보(소득, 상환 부담, 최근 계좌 개설)가 남아 있었다.
\item \textbf{평가 표본 자체가 편향될 수 있다.} 이미 끝난 대출만 보면 부도확률 순위가 기대수익 순위를 이기는 것처럼 보인다(1--2년 관측 시 $-3.83$\%p). 전체를 보면 반대(+0.91\%p)다.
\end{enumerate}
\Intuit{좋은 모델보다 정직한 성적표가 먼저다. 무엇을 재고, 언제의 정보를 쓰고, 어떤 표본으로 채점하는지가 결론을 좌우했다.}
\end{keybox}

\begin{warnbox}[처음 판단이 틀렸던 것(정정 기록에서)]
\begin{itemize}[leftmargin=*]
\item 지난 우수 보고서의 ``부도율 24\%, 33\%''는 부도 \ensuremath{\div} 정상 비율이었다. 실제 부도율은 19.3\%, 25.1\%.
\item ``부도확률 순위는 무작위와 비슷하다''(시드 1개)는 틀렸다. 시드 5개와 부트스트랩으로 보면 유의하게 양수이고 해마다 커진다.
\item ``선별 곡선은 새로운 방법이다''는 틀렸다. 수익 기반 점수(Serrano-Cinca \& Guti\'errez-Nieto, 2016), EMP(Verbraken 외, 2014)가 이미 있다.
\item 로지스틱이 2016년 AUC 0.508로 무너진 원인은 결측 처리가 아니라 \textbf{학습 때 없던 변수}였다.
\item ``CatBoost는 새 컬럼을 조기 종결의 대리 변수로 배웠다''는 가설은 반증되었다(부도율 17.2\% vs 17.6\%).
\item 2017년 보류 구간의 ``부도확률 순위 우위''는 종결 편향으로 재현되었다.
\end{itemize}
\end{warnbox}

\chapter{과제와 데이터}

\section{과제가 실제로 묻는 것}
과제 문서는 ``Sharpe ratio를 최대로 만드는 대출 승인 전략''을 요구한다. 조교 답변에 따르면 Sharpe는 바꾸거나 다른 식을 써도 되고, accuracy는 중요하지 않으며, \textbf{과정과 논리 도출}이 평가의 중심이다. 그래서 질문을 이렇게 다시 쓴다.

\begin{keybox}[이 프로젝트의 질문]
투자자가 LendingClub(LC) 대출 중 무엇을 사야 하는가? 그리고 그 판단이 ``운이 아니라 정보 때문''이라는 것을 어떻게 보일 수 있는가?
\Intuit{예측을 잘하는 대회가 아니라, 돈이 되는 고르기를 정직하게 증명하는 숙제다.}
\end{keybox}

질문이 두 부분이라는 점이 중요하다. 앞부분(무엇을 사나)은 모델과 규칙의 문제이고, 뒷부분(어떻게 보이나)은 평가와 검증의 문제다. 이 노트의 절반 이상이 뒷부분에 쓰이는 이유다.

\section{LendingClub이라는 시장}
LC는 2007년부터 2020년까지 운영된 미국의 P2P 대출 플랫폼이다. 돈을 빌리려는 사람이 신청하면 LC가 심사해 등급(A--G, 세부등급 A1--G5 35개)과 금리를 붙이고, 투자자는 그 대출 조각(note)을 산다. 대출은 무담보 개인 신용대출이고 만기는 36개월 또는 60개월이다. 2021년 2월 Radius Bank 인수를 마치며 은행이 되었고 P2P 투자 상품은 접었다.

\begin{insight}[투자자 입장에서의 세 가지 위험]
\begin{enumerate}[leftmargin=*]
\item \textbf{부도(default)}: 갚다가 멈춘다. 원금 일부를 잃는다.
\item \textbf{조기상환(prepayment)}: 일찍 다 갚는다. 원금은 돌아오지만 높은 이자를 받을 기간이 줄어든다.
\item \textbf{재투자}: 매달 돌아오는 원리금을 다시 어딘가에 넣어야 한다. 그 수익률이 전체 성과를 바꾼다.
\end{enumerate}
\end{insight}

LC가 매긴 금리는 이미 부도 위험의 가격이다. 그래서 ``부도가 적은 대출 고르기''는 정답이 아니다. 금리가 낮은 A등급만 고르면 부도는 적지만 버는 것도 적다. 이 긴장이 프로젝트 전체를 관통한다.

\section{데이터 한눈에}
\begin{table}[H]\centering\small
\begin{tabularx}{\linewidth}{@{}lL@{}}\toprule
항목 & 값\\\midrule
공식 train & 1,755,294건 \ensuremath{\times} 141개 컬럼, 2007--2020년 발행\\
결과 상태 & Fully Paid 약 90만, Current(진행 중) 약 62만, Charged Off 약 22만, 기타 소수\\
종결 표본 & 결과가 확정된 대출 1,116,156건(LC 신용정책 미충족 대출 제외)\\
모델 특성 & 승인 시점에 알 수 있는 기본 101개(\ft{src/lc/data.py}의 \ft{feature_columns})\\
주 평가 표본 & 2013--2016년 발행 36개월 종결 대출, 발행월 48개\\
공식 test & 조교 답변: 일주일 안에 배포, 상환 결과 포함, 진행 중 대출은 없을 것\\\bottomrule
\end{tabularx}
\caption{데이터 요약.}
\end{table}

\section{컬럼을 시점으로 나누기: Data leakage}
141개 컬럼을 보면 가장 먼저 던질 질문은 하나다.

\begin{keybox}[누수 판별 질문]
``대출을 승인하던 날, 이 값을 알 수 있었나?''
알 수 없었다면 그 컬럼은 학습에 쓰면 안 된다. 쓰면 \textbf{data leakage}다.
\Intuit{시험 답안지를 보고 푼 모의고사 점수는 실력이 아니다.}
\end{keybox}

\begin{table}[H]\centering\small
\begin{tabularx}{\linewidth}{@{}p{33mm}LL@{}}\toprule
묶음 & 예 & 처리\\\midrule
승인 시점 정보 & \ft{loan_amnt}, \ft{int_rate}, \ft{grade}, \ft{annual_inc}, \ft{dti}, \ft{fico_range_low}, \ft{emp_length}, \ft{addr_state} & 특성으로 사용\\
상환 결과 & \ft{total_pymnt}, \ft{recoveries}, \ft{last_pymnt_d}, \ft{out_prncp} & 특성 금지, \textbf{수익 계산에만} 사용\\
사후 갱신 & \ft{last_fico_range_high}, \ft{hardship_*}, \ft{settlement_*}, \ft{debt_settlement_flag} & 삭제(\ft{LEAK}, \ft{LEAK_PREFIX})\\
식별, 설명 & \ft{id}, \ft{url}, \ft{desc}, \ft{title} & 삭제\\\bottomrule
\end{tabularx}
\caption{컬럼 묶음. 실제 목록은 \ft{src/lc/data.py}.}
\end{table}

\begin{warnbox}[미묘한 누수]
\ft{last_fico_range_high}는 이름만 보면 신용점수지만, 대출 후 마지막으로 갱신된 점수다. 부도가 나면 점수가 떨어지므로 결과를 거의 그대로 알려 준다. 이런 컬럼 하나가 AUC를 0.9 넘게 올리는데, 그 점수는 아무 의미가 없다.
\end{warnbox}

\section{결과 라벨 만들기}
\ft{data.parse}는 원본에 다음 열을 붙인다.
\begin{itemize}[leftmargin=*]
\item \ft{outcome}: Fully Paid, Charged Off 등 상태를 정리한 값. \ft{resolved}: 결과가 확정되었는가.
\item \ft{bad}: 부도면 1. Charged Off와 Default를 부도로 본다.
\item \ft{r}: 약정 금리(소수), \ft{term_m}: 만기 개월, \ft{tau}: 발행부터 마지막 납부까지 개월 수.
\item \ft{policy_fail}: ``Does not meet the credit policy''로 표시된 옛 대출. 심사 기준이 달라 분석에서 뺀다.
\end{itemize}

\begin{career}
``누수를 어떻게 막았나요?''는 면접 단골 질문이다. 답의 핵심은 목록이 아니라 \textbf{판별 기준}이다. ``각 컬럼에 대해 결정 시점에 관측 가능한지를 물었고, 결과 컬럼은 라벨과 수익 계산에만 쓰고 특성에서는 뺐다''라고 말하면 된다. 이 프로젝트에서는 한 단계 더 나아가, 컬럼뿐 아니라 \textbf{행(어떤 대출을 학습에 넣나)}도 시점을 지켰다(3장).
\end{career}

\chapter{수익률을 어떻게 잴 것인가}

과제가 Sharpe를 말하는 순간 ``대출 하나의 수익률''이 필요해진다. 그런데 데이터에는 수익률 컬럼이 없다. 있는 것은 약정 금리, 월 상환액, 받은 총액, 마지막 납부일뿐이다. 이 장은 그것들로 수익률을 만드는 법과, 정의를 어떻게 고르느냐에 따라 결론이 뒤집히는 이유를 다룬다.

\section{원리금 균등 상환}
\textbf{질문.} 1만 달러를 연 12\%로 36개월 동안 빌리면 매달 얼마를 갚나?

매달 같은 금액 $A$를 내는데, 그 36번의 납입을 오늘 가치로 할인해 더하면 빌린 돈 $P$가 되어야 한다. 월 이자율 $i=r/12$로 할인하면
\[
P=\sum_{t=1}^{T}\frac{A}{(1+i)^t}=A\cdot\frac{1-(1+i)^{-T}}{i}
\quad\Longrightarrow\quad
A=P\,\frac{i}{1-(1+i)^{-T}}.
\]

\begin{workout}[월 상환액]
\Use{원리금 균등 상환식(이 절), 등비수열의 합}
\Step{1} 월 이자율: $i=0.12/12=0.01$. \Concept{연 금리를 월로}
\Step{2} 할인 인자: $(1.01)^{-36}=0.6989$. 따라서 $1-(1.01)^{-36}=0.3011$.
\Step{3} $A=10000\times 0.01/0.3011=332.14$. \Concept{원리금 균등 상환}
\Ans{매달 \$332.14. 데이터의 \ft{installment}가 바로 이 값이다.}
\Verify{$332.14\times 36=11{,}957$. 이자 총액 약 \$1,957로, 원금이 줄어드는 만큼 12\% \ensuremath{\times} 3년(\$3,600)보다 작다.}
\Trap{연 금리 12\%를 그대로 36으로 곱해 이자를 계산한다. 원금이 매달 줄어드는 것을 잊은 것이다.}
\end{workout}

\section{IRR: 현금흐름이 실제로 낸 수익률}
\textbf{질문.} 대출이 중간에 부도 나면 수익률은 몇인가?

받은 돈의 흐름을 $CF_0=-P$(투자), $CF_1,\dots,CF_H$(받은 돈)로 적자. 이 흐름의 현재가치를 0으로 만드는 월 할인율 $m$이 있으면, 그것이 이 투자가 실제로 낸 월 수익률이다.
\[
\sum_{t=0}^{H}\frac{CF_t}{(1+m)^t}=0,\qquad \mathrm{IRR}=(1+m)^{12}-1.
\]
식이 다항식이라 닫힌 해가 없으므로 수치로 푼다. 코드는 \ft{returns.irr}이며 Brent 방법(\ft{brentq})으로 근을 찾는다.

\begin{keybox}[IRR]
돈을 넣고 돌려받은 흐름 전체를 하나의 ``연 이자율''로 요약한 값이다. 정상 상환이면 약정 금리와 같고, 부도 나면 그보다 낮다.
\Intuit{이 통장에 넣었다 뺀 돈의 흐름을, 매년 몇 \% 이자 통장과 같다고 볼 수 있나?}
\end{keybox}

\begin{workout}[정상 상환 대출의 IRR]
\Use{IRR 정의(이 절), 월 복리}
\Step{1} 현금흐름: $CF_0=-10000$, $CF_1=\dots=CF_{36}=332.14$.
\Step{2} $m=0.01$을 넣으면 \S2.1의 식에 의해 현재가치가 정확히 0이다. 따라서 월 IRR은 1\%. \Concept{IRR 정의}
\Step{3} 연환산: $(1.01)^{12}-1=0.1268$. \Concept{월 복리}
\Ans{연 IRR 12.68\%. 약정 금리 12\%는 단리식 표기(APR)이고, IRR은 복리라 조금 크다.}
\Verify{\ft{checks.py} 2번: 월 0.01000, 연 0.12683.}
\end{workout}

\begin{workout}[10개월 내고 부도 난 대출]
\Use{현금흐름 복원(아래), IRR, 손실률}
\Step{1} 상황(예시): 10번 냈고, 15개월째 회수금 \$500이 들어왔다. 받은 총액 $=10\times332.14+500=3{,}821.43$.
\Step{2} 손실률 $=1-3821.43/10000=0.618$. \Concept{LGD 정의, \S5.6}
\Step{3} 현금흐름 $CF_0=-10000$, $CF_{1..10}=332.14$, $CF_{15}=500$의 IRR을 수치로 풀면 연 $-76.8\%$.
\Step{4} 같은 대출의 단순 연환산(우승 논문식): $(3821.43/10000)^{12/36}-1=-27.4\%$.
\Ans{IRR $-76.8\%$, 단순 연환산 $-27.4\%$. 같은 대출인데 정의에 따라 세 배 가까이 다르다.}
\Verify{\ft{checks.py} 3번. 손실률 0.618은 우연히 2013년 학습 표본의 실제 평균 LGD(0.618)와 같다.}
\Trap{부도 대출의 IRR이 $-100\%$에 가까울 것이라 생각한다. 이미 받은 원리금이 있어 그보다 덜 나쁘다. 반대로 단순 연환산은 손실을 3년에 걸쳐 나눠 덜 나빠 보이게 만든다.}
\end{workout}

\paragraph{현금흐름 복원.} 데이터에는 월별 납입 기록이 없다. 그래서 \ft{returns.schedule}은 이렇게 근사한다: 받은 총액을 월 상환액으로 나눈 몫 $k_f$만큼은 매달 정상 납입했고, 나머지 $\text{rem}=\text{받은 총액}-k_f A$는 마지막 납부 달($\tau$)에 한꺼번에 받았다. 조기상환은 큰 rem이 일찍 들어오는 것으로, 부도 후 회수는 작은 rem으로 표현된다. 이 근사는 연구 노트의 가정 목록에 올라 있다.

\section{포트폴리오 IRR과 수수료}
투자자는 대출 하나가 아니라 여러 개를 산다. 같은 달에 발행된 대출들의 현금흐름을 더해 한 포트폴리오로 보고 IRR을 구한다(\ft{returns.pooled_irr}). LC는 투자자가 받은 돈에서 1\% 서비스 수수료를 뗀다.

\begin{workout}[정상 9건 + 부도 1건]
\Use{pooled IRR(이 절), 수수료}
\Step{1} 각 \$10,000, 연 12\%, 36개월. 9건은 정상 상환, 1건은 앞 예제처럼 부도.
\Step{2} 월별 합계: 1--10개월 $10\times332.14$, 11--36개월 $9\times332.14$, 15개월에 +500. 0개월 $-100{,}000$.
\Step{3} 받은 돈에 0.99를 곱한다. \Concept{투자자 수수료 1\%}
\Step{4} IRR을 풀면 연 6.85\%.
\Ans{부도율 10\%가 연 12.68\%의 수익을 6.85\%로 깎았다.}
\Verify{\ft{checks.py} 4번.}
\end{workout}

\section{네 가지 수익률 정의와 결론의 뒤집힘}
\begin{table}[H]\centering\small
\begin{tabularx}{\linewidth}{@{}p{34mm}LL@{}}\toprule
정의 & 계산 & 문제점\\\midrule
수업식 & 정상이면 약정 금리, 부도면 0 & 부도 손실(원금)을 무시해 과대평가\\
단순 연환산 & (받은/대출)$^{12/T}-1$ & 돈이 일찍 돌아온 것(재투자)을 무시, 조기상환을 손해로 셈\\
재투자 & 받은 돈을 국채 금리로 만기까지 굴림 & 재투자 금리 가정 필요\\
코호트 IRR & 발행월 포트폴리오의 현금흐름 IRR, 수수료 1\% & 현금흐름 시점이 근사\\\bottomrule
\end{tabularx}
\caption{수익률 정의 네 가지(\ft{src/lc/returns.py}). 이 프로젝트의 주 정의는 마지막 줄이다.}
\end{table}

\begin{keybox}[결과 C1--C3]
2013--16년 36개월 대출을 \textbf{전부 승인}하면 코호트 IRR 5.73\%로 같은 기간 3년 국채(0.87\%)를 48개 발행월 모두에서 이긴다(C2). 지난 우승 논문은 ``전부 승인 \ensuremath{<} 국채''라고 결론냈는데, 이는 단순 연환산과 최근 발행분의 종결 표본에서 비롯된 것으로 보인다(C3, 잠정).
\Intuit{같은 데이터라도 자를 대는 방법이 다르면 다른 답이 나온다. 그래서 자부터 정해야 한다.}
\end{keybox}

\section{비교 기준: 무위험수익률은 무엇인가}
Sharpe와 초과수익에는 ``아무 위험 없이 벌 수 있는 수익'' $r_f$가 필요하다.
\begin{itemize}[leftmargin=*]
\item \textbf{만기를 맞춘 국채}(FRED GS3, GS5): 36개월 대출에는 3년 국채, 60개월에는 5년 국채를 발행월 기준으로 붙인다(\ft{returns.attach_treasury}). 투자 기간이 같은 대안이라 가장 정직하다.
\item \textbf{3개월 국채}(TB3MS): 교과서의 $r_f$. 만기가 달라 기간 프리미엄이 섞인다.
\item \textbf{회사채 BAA}: 무위험은 아니지만 ``비슷한 위험의 다른 투자''라는 실무적 비교 기준. BAA로 바꾸면 전부 승인의 우위가 사라지고 선별의 가치가 커진다(C8).
\end{itemize}

\begin{insight}
실무에서 대출 포트폴리오의 성과는 조달 비용(은행이면 예금 금리 + 스프레드) 또는 같은 위험의 벤치마크 대비로 본다. ``무위험''은 출발점일 뿐이라, 이 프로젝트는 국채를 주 기준으로 두고 BAA와 +1--3\%p 프리미엄을 민감도로 함께 보고한다.
\end{insight}

\begin{career}
``Accuracy가 아니라 무엇을 최적화했나?''라는 질문에 이 장이 답이다. 라벨(부도 0/1)이 아니라 \textbf{현금흐름에서 복원한 IRR}을 목표로 삼았고, 정의 선택이 결론을 뒤집는다는 것을 수치로 보였다.
\end{career}

\chapter{시간을 지키는 법: 시점 학습과 편향}

2장에서 수익률의 자를 정했다. 이 장은 그 자로 \textbf{언제의 정보로} 학습하고 \textbf{어떤 대출로} 채점하는지를 다룬다. 누수는 컬럼(1장)에서만 생기지 않는다. 행에서도 생긴다.

\section{무작위 분할은 왜 성적을 부풀리나}
보통 머신러닝 수업에서는 데이터를 무작위로 80:20으로 나눈다. LC 데이터에 그렇게 하면 2016년 대출로 학습해 2013년 대출을 맞히는 일이 생긴다. 2016년의 경기, LC의 심사 기준, 금리 수준은 2013년 투자자가 알 수 없던 정보다.

\begin{keybox}[결과 C5]
무작위 분할 AUC 약 0.73, 시점 학습 AUC 0.62--0.68. 차이 전부가 ``미래를 본 값''이다.
\Intuit{내년 기출문제로 공부하고 올해 시험을 보면 점수는 높지만, 실제로는 그런 일이 불가능하다.}
\end{keybox}

\section{시점 학습(as-of training)}
연도 $Y$에 발행된 대출을 평가할 때, 모델은 \textbf{$Y$년 1월 1일에 이미 결과가 확정된 대출}로만 학습한다(\ft{asof.training_set}). 확정 조건은 두 가지다.

\begin{enumerate}[leftmargin=*]
\item 그 날짜 전에 발행되었다.
\item 그 날짜까지 결과가 났다. 정상 상환은 마지막 납부 달에, 부도는 \textbf{마지막 납부 + 5개월}에 확정된 것으로 본다(\ft{CHARGE_OFF_LAG = 5}).
\end{enumerate}

\begin{figure}[H]\centering
\begin{tikzpicture}[x=1.0cm,y=0.75cm,font=\small]
\draw[-{Stealth}] (0,0)--(9.6,0) node[right]{시간};
\foreach \x/\l in {1/2010,3/2011,5/2012,7/2013,9/2014} {\draw (\x,0.1)--(\x,-0.1) node[below]{\l};}
\draw[very thick] (7,-0.6)--(7,4.3) node[above]{기준일 2013-01-01};
\draw[black!70,line width=2pt] (1.2,1)--(4.5,1); \node[right] at (9.8,1) {A: 정상 상환, 확정};
\draw[black!70,line width=2pt] (2,2)--(5.4,2); \draw[black!35,line width=2pt,dashed] (5.4,2)--(6.4,2); \node[right] at (9.8,2) {B: 부도, 5개월 뒤 확정};
\draw[black!70,line width=2pt] (4,3)--(6.6,3); \draw[black!35,line width=2pt,dashed] (6.6,3)--(7.6,3); \node[right] at (9.8,3) {C: 기준일 뒤 확정, 제외};
\draw[black!70,line width=2pt] (5.5,4)--(9.5,4); \node[right] at (9.8,4) {D: 진행 중, 제외};
\end{tikzpicture}
\caption{시점 학습의 학습 표본. 2013년 대출을 평가할 모델은 A와 B로만 학습한다. 점선은 상각 지연 5개월.}
\end{figure}

\begin{whybox}[왜 5개월인가]
LC는 연체 120일(약 4개월)이 넘으면 상각(charge-off)한다. 마지막 납부 이후 상각이 장부에 잡히기까지 약 5개월이 걸린다. 지연을 0으로 두면 아직 ``연체 중''이던 대출을 이미 부도로 알았던 것처럼 학습해 미래 정보가 섞인다. 6단계에서 지연 0과 5를 비교했고, 가정에 따라 편향 방향이 바뀌는 것을 확인했다.
\end{whybox}

\paragraph{대가.} 시점 학습은 표본이 작다. 2013년 모델의 학습 대출은 11,345건이고, 2016년 모델은 156,153건이다. 이른 해일수록 모델이 약하고 불안정하다(5장의 2013년 결과들이 흔들리는 이유).

\section{종결 표본 편향(resolution bias)}
\textbf{질문.} 평가할 때 ``이미 결과가 난 대출''만 쓰면 무엇이 달라지나?

\begin{workout}[4건짜리 종결 편향]
\Use{censoring, 종결 표본 편향(이 절)}
\Step{1} 같은 달 발행된 36개월 대출 4건(예시):
A 저금리 6\%, 3년 내내 정상. B 고금리 20\%, 3년 내내 정상. C 고금리 20\%, 6개월 뒤 부도. D 저금리 6\%, 1년 안에 조기상환.
\Step{2} 3년 뒤 전체를 보면 부도율 $1/4=25\%$. 고금리 B의 이자가 C의 손실을 메운다.
\Step{3} 발행 1년 뒤에 ``끝난 대출만'' 보면 C와 D만 남는다. 부도율 $1/2=50\%$. \Concept{censoring}
\Step{4} 남은 표본에서 ``고금리를 고르는 규칙''은 C만 고르게 되어 최악이고, ``부도확률이 낮은 것을 고르는 규칙''은 D를 골라 최선이다.
\Ans{일찍 본 표본은 부도를 과대 대표하고, 오래 이자를 받아야 이득인 고금리 정상 대출(B)을 숨긴다.}
\Trap{``종결 대출만 쓰면 깨끗하다''고 생각한다. 깨끗한 것은 라벨이고, 표본 구성은 오히려 왜곡된다.}
\end{workout}

\subsection*{실제 데이터로 시뮬레이션(26단계)}
2013--15년 36개월 대출은 지금 모두 끝나 정답을 안다. 같은 대출을 관측 시점에 맞춰 일부러 잘라, 두 순위 방식의 선별 가치(6장)를 비교했다. LightGBM 시드 10개, 발행월 부트스트랩 1만 번, 판정 기준은 실행 전에 고정했다.

\begin{table}[H]\centering\small
\begin{tabular}{@{}lrrrrr@{}}\toprule
관측 & 남은 비율 & 부도율 & 기대수익 순위 & 부도확률 순위 & 차이(95\% 구간)\\\midrule
발행 후 1--2년 & 약 30\% & 약 18\% & 1.44\%p & 5.28\%p & $-3.83$ ($-4.67$, $-3.04$)\\
발행 후 2--3년 & 약 59\% & 약 18\% & 1.18\%p & 2.21\%p & $-1.03$ ($-1.36$, $-0.70$)\\
전체 & 100\% & 약 14\% & 1.01\%p & 0.10\%p & $+0.91$ ($+0.73$, $+1.10$)\\\bottomrule
\end{tabular}
\caption{종결 편향 시뮬레이션(\ft{scripts/26_resolution_bias.py}). 단위는 무작위 승인 대비 \%p.}
\end{table}

\begin{keybox}[결과 C23과 C6 정정]
관측 기간이 짧을수록 부도확률 순위가 이기는 것처럼 보이고, 전체를 보면 기대수익 순위가 이긴다. 14단계의 2017년 보류 구간에서 부도확률 순위가 이긴 것은 이 편향으로 설명된다.
\Intuit{일찍 채점하면 ``빨리 망한 대출''만 보이므로, 망하지 않는 것만 신경 쓰는 규칙이 똑똑해 보인다.}
\end{keybox}

\begin{warnbox}[고칠 수 있나]
종결 여부는 결과(부도, 조기상환) 자체에 달려 있다. 이런 결측은 MNAR(4장)이라, 특성으로 만든 가중치(inverse probability weighting)로는 고칠 수 없다. 대신 \textbf{test를 보기 전에 읽는 법을 정했다}: 발행연도별 만기 도달 비율부터 확인하고, 비율이 높은 연도만 두 순위 비교의 근거로 쓴다.
\end{warnbox}

\section{생존분석은 왜 안 썼나}
진행 중 대출 62만 건을 버리는 것이 아깝다면, 생존분석(survival analysis)으로 ``아직 안 끝난'' 정보까지 쓸 수 있다. 시도했지만 단순 PD 모델을 이기지 못했다(연구 노트 실패 목록). 수준(확률 크기)에는 편향이 있어도 순위는 크게 무너지지 않았기 때문이다.

\begin{career}
시계열이 섞인 데이터에서 무작위 분할을 쓰지 않은 이유, 그리고 \textbf{평가 표본의 선택 편향}까지 시뮬레이션으로 정량화한 것은 강한 이야깃거리다. ``모델 성능을 올렸다''보다 ``성능이 부풀려지는 경로 두 가지(무작위 분할, 종결 편향)를 찾아 크기를 쟀다''가 훨씬 드물다.
\end{career}

\chapter{EDA, 피처, 결측}

\section{EDA는 무엇을 위한 것인가}
EDA(exploratory data analysis)는 그림을 많이 그리는 일이 아니라, \textbf{모델링 결정을 내리기 위한 질문에 답하는 일}이다. 이 프로젝트의 EDA 질문은 네 개였다.
\begin{enumerate}[leftmargin=*]
\item 어떤 변수가 혼자서 부도를 가르는가? \ensuremath{\to} IV(\S4.2)
\item 어떤 변수들이 서로 같은 정보인가? \ensuremath{\to} 순위상관
\item LC 등급이 이미 설명한 것 말고 남은 정보가 있는가? \ensuremath{\to} 등급 안 분석(\S4.3)
\item 값이 왜 비어 있나? \ensuremath{\to} 결측 메커니즘(\S4.4)
\end{enumerate}
12단계 EDA는 평가 구간을 보지 않도록 2007--12년 대출만 썼다. 20단계 EDA(\ft{results/eda/})는 2007--16년 종결 대출을 썼다.

\section{WOE와 IV: 한 변수의 힘}
\textbf{질문.} 연소득이 부도와 얼마나 관련 있는지를 한 숫자로 말할 수 있을까?

변수를 구간(bin)으로 나누고, 각 구간에 정상 대출과 부도 대출이 얼마나 몰렸는지 비교한다. 전체 정상 $G$건 중 구간 $k$에 $g_k$건, 전체 부도 $B$건 중 $b_k$건이 있으면
\[
\mathrm{WOE}_k=\ln\frac{g_k/G}{b_k/B},\qquad
\mathrm{IV}=\sum_k\Big(\frac{g_k}{G}-\frac{b_k}{B}\Big)\mathrm{WOE}_k .
\]
WOE가 양수면 그 구간은 정상 쪽, 음수면 부도 쪽 증거다. IV는 모든 구간의 ``치우침 \ensuremath{\times} 증거''를 더한 것이라 항상 0 이상이다. 관행상 0.02 미만은 약함, 0.1--0.3 중간, 0.3 이상 강함이다.

\begin{workout}[결측 구간의 WOE]
\Use{WOE, IV(이 절), MNAR(\S4.4)}
\Step{1} 예시: 정상 900건, 부도 100건. 어떤 변수가 결측인 구간과 값이 있는 구간 두 개.
\Step{2} 경우 1: 결측 구간에 정상 270, 부도 30. 비중이 정상 $270/900=0.3$, 부도 $30/100=0.3$로 같다.
$\mathrm{WOE}=\ln(0.3/0.3)=0$, 값 있는 구간도 0. IV $=0$. \Concept{WOE 정의}
\Step{3} 경우 2: 결측 구간에 정상 240, 부도 60. 비중 $0.267$ vs $0.6$.
$\mathrm{WOE}_{\text{결측}}=\ln(0.267/0.6)=-0.811$, $\mathrm{WOE}_{\text{있음}}=\ln(0.733/0.4)=+0.606$.
\Step{4} $\mathrm{IV}=(0.267-0.6)(-0.811)+(0.733-0.4)(0.606)=0.472$. \Concept{IV 정의}
\Ans{경우 1은 결측이 아무 정보도 없고(IV 0), 경우 2는 ``비어 있다는 사실'' 자체가 강한 부도 신호다(IV 0.47).}
\Verify{\ft{checks.py} 8번.}
\Trap{결측을 평균으로 채워 버린다. 경우 2라면 가장 강한 신호를 지운 것이다.}
\end{workout}

\begin{table}[H]\centering\small
\begin{tabular}{@{}lr|lr@{}}\toprule
변수 & IV & 변수 & IV\\\midrule
\ft{sub_grade} & 0.487 & \ft{term} & 0.178\\
\ft{grade} & 0.452 & FICO & 0.112\\
\ft{int_rate} & 0.441 & \ft{annual_inc} & 0.029\\\bottomrule
\end{tabular}
\caption{IV 상위(20단계). 101개 중 67개는 혼자서는 IV 0.02 미만이다.}
\end{table}

등급, 세부등급, 금리가 압도적이다. 셋은 사실상 같은 정보(LC의 판단)다. 순위상관 0.9 이상인 쌍이 12개 있었다(FICO 상한과 하한 1.0, 대출액과 승인액 1.0 등).

\section{등급 안에 남은 정보}
IV만 보면 연소득은 약한 변수(0.029)다. 그런데 \textbf{같은 세부등급 안에서} 비교하면 이야기가 달라진다.

\begin{keybox}[결과 C14, C18]
등급이 설명하지 못한 정보는 소득, 상환 부담(월 상환액/월 소득), 최근 계좌 개설, 부채비율(DTI), 지역에 남아 있다. 같은 세부등급, 같은 금리에서 소득 상위이면서 DTI 하위 20\%인 대출이 IRR이 약 0.9\%p 높다(모델 없이 계산).
\Intuit{LC는 신용점수는 가격에 잘 넣었지만, ``갚을 능력''은 덜 넣었다.}
\end{keybox}

혼자서는 약하고 등급과 함께 볼 때 힘이 나는 변수가 있다는 것, 이것이 후보 A(같은 등급 안 선별)의 출발점이다. 단, 소득 효과는 해마다 줄었다(2.66\%p에서 0.19\%p). 시장이 그 정보를 점점 가격에 넣었을 수 있다.

\section{결측 메커니즘}
\textbf{질문.} 값이 비어 있는 이유가 결과와 관련 있나?

\begin{table}[H]\centering\small
\begin{tabularx}{\linewidth}{@{}p{30mm}LL@{}}\toprule
메커니즘 & 뜻 & LC 예\\\midrule
MCAR & 완전히 무작위로 비어 있음 & 거의 없음(소량 결측 일부)\\
MAR & 관측된 다른 변수(예: 발행 시기)로 설명됨 & 2015년 말 이후 생긴 새 변수 14개(\ft{il_util} 등)\\
MNAR & 비어 있는 값 자체와 관련 있음 & \ft{emp_length}: 무직이거나 밝히기 꺼리면 비움\\
구조적 & 정의상 값이 없음 & 공동신청 변수 15개(단독 신청이면 없음)\\
사건 없음 & 그 사건이 없었음 & \ft{mths_since_last_delinq}: 연체한 적이 없으면 빔\\\bottomrule
\end{tabularx}
\caption{결측 유형(\ft{src/lc/missing.py}의 \ft{classify}).}
\end{table}

\begin{whybox}[사건 없음은 ``아주 오래전''이다]
``마지막 연체 후 몇 개월''이 비었다는 것은 연체한 적이 없다는 뜻이다. 평균(예: 30개월)으로 채우면 ``30개월 전에 연체했다''는 거짓 정보가 된다. 그래서 선형 모델용으로는 학습 데이터의 최댓값으로 채우고 결측 표시를 붙이는 방식을 시험했다(\ft{typed_linear}).
\end{whybox}

\section{결측 처리 결정(22단계)}
결측 처리는 모델 비교보다 \textbf{먼저} 정해야 한다. 그렇지 않으면 ``모델 A가 낫다''가 사실은 ``모델 A에 맞춘 결측 처리가 낫다''일 수 있다. 판정 기준은 실험 전에 정했다: 가장 단순한 방식보다 선별 가치 0.1\%p 초과 또는 AUC 0.002 초과로 나은 연도가 4개 중 3개 이상이면 채택.

\begin{itemize}[leftmargin=*]
\item \textbf{트리 모델}(LightGBM): 결측 그대로 / 중앙값 / 중앙값 + 결측 표시 / 유형별 처리 모두 시드 차이 안. \ensuremath{\to} \textbf{결측 그대로}. LightGBM은 분할마다 결측을 왼쪽, 오른쪽 중 나은 쪽으로 보내는 법을 스스로 배운다.
\item \textbf{선형 모델}(로지스틱): 1차 비교에서 2016년 AUC 0.508로 무너졌다.
\end{itemize}

\begin{warnbox}[로지스틱 붕괴의 진짜 원인]
결측 처리가 문제가 아니었다. 2016년 모델의 학습 기간(2015년까지)에는 새 변수 33개가 거의 비어 있었다. 학습에서 그 변수의 표준편차는 소수 대출로 계산되어 매우 작았고, 평가 연도에는 모든 대출에 값이 있어 표준화 값이 최대 290까지 튀었다. 선형 모델은 큰 입력에 계수를 그대로 곱하므로 예측이 폭발한다.
\end{warnbox}

\begin{keybox}[확정 방침]
\begin{itemize}[leftmargin=*]
\item 트리: 결측 그대로.
\item 선형: \textbf{가용성 규칙}(학습 가용률 50\% 미만 변수 제외) \ensuremath{\to} 1--99\% 분위 자르기 \ensuremath{\to} 중앙값 대치. \ft{emp_length} 결측은 ``모름'' 범주로.
\end{itemize}
두 규칙을 넣자 2016년 AUC가 0.678로 정상화됐다. 이후 비교에서는 중앙값보다 확실히 나은 방식이 없었다.
\Intuit{결측을 어떻게 채우느냐보다, 학습 때 거의 없던 변수를 쓰지 않는 것이 훨씬 중요했다.}
\end{keybox}

\section{Feature engineering}
만든 피처 6종: 상환 부담(PTI = 월 상환액/월 소득), 대출 대비 소득(LTI), 신용 이력 개월, 근속 연수 숫자화, 우편번호 3자리 부도율(out-of-fold target encoding), FICO 평균. 대조군으로 노이즈 피처와 시드 변경을 함께 넣었다.

\begin{keybox}[결과 C15, C16, C21]
LightGBM에서는 모든 피처의 효과($-0.01$--$+0.02$\%p)가 시드만 바꿔도 생기는 차이(최대 0.09\%p)보다 작았다. 트리는 원본 변수(소득, 상환액, 대출액)에서 이미 비율 관계를 잡는다. 선형 모델에서는 피처가 조금 도움이 된다(L2 \ensuremath{\to} L2 + 피처, Lasso \ensuremath{\to} Lasso + 피처).
\Intuit{나무는 나눗셈을 스스로 흉내 내지만, 직선 모델은 누가 나눠 줘야 한다.}
\end{keybox}

\begin{prof}[target encoding의 누수]
우편번호별 부도율을 그 대출 자신이 포함된 데이터로 계산하면 자기 라벨이 특성에 섞인다. 학습 데이터를 5조각으로 나눠 각 조각의 값은 나머지 4조각으로 계산(out-of-fold)하고, 표본이 적은 우편번호는 전체 평균 쪽으로 당겼다(smoothing).
\end{prof}

\begin{career}
``Feature engineering으로 성능을 얼마나 올렸나?''에 ``0이었다, 그리고 그걸 어떻게 확인했나''로 답할 수 있다. 노이즈 피처와 시드 변경을 \textbf{대조군}으로 넣어 효과 크기를 잡음과 비교한 것이 핵심이다.
\end{career}

\chapter{모델: 부도확률에서 기대수익까지}

모델이 하는 일은 하나다. 대출의 특성 $x$를 받아 부도확률 $p=P(\text{부도}\mid x)$를 내는 것. 그 $p$를 금리, 손실률, 국채와 합쳐 \textbf{기대수익 점수}를 만들고(\S5.6), 점수로 승인한다.

\section{Logistic regression}
\textbf{질문.} 특성의 가중합 $z=\beta_0+\beta^\top x$는 $-\infty$에서 $+\infty$까지 아무 값이나 된다. 이걸 0과 1 사이 확률로 바꾸려면?

조건은 세 가지다: $z$가 클수록 확률이 커지고, $z\to-\infty$면 0, $z\to+\infty$면 1. 이를 만족하는 가장 단순한 S자 함수가 sigmoid다.
\[
p=\sigma(z)=\frac{1}{1+e^{-z}},\qquad \ln\frac{p}{1-p}=z .
\]
오른쪽 식이 핵심이다. \textbf{log-odds}가 특성의 선형식이다. 계수 $\beta_j$는 ``$x_j$가 1 늘 때 log-odds가 $\beta_j$ 는다''로 읽는다.

\begin{workout}[sigmoid에 숫자 넣기]
\Use{sigmoid, log-odds(이 절)}
\Step{1} 예시: $z=-2+0.5x$, $x$는 ``최근 6개월 신규 계좌 수''라 하자.
\Step{2} $x=0$: $z=-2$, $p=1/(1+e^{2})=0.119$. \Concept{sigmoid}
\Step{3} $x=2$: $z=-1$, $p=1/(1+e^{1})=0.269$.
\Step{4} $x=4$: $z=0$, $p=0.5$.
\Ans{계좌 2개 늘 때마다 odds가 $e^{1}=2.72$배가 된다. 확률 증가폭은 일정하지 않다(0.12 \ensuremath{\to} 0.27 \ensuremath{\to} 0.50).}
\Verify{\ft{checks.py} 13번.}
\end{workout}

학습은 \textbf{log loss}(교차 엔트로피) $-\frac1n\sum[y_i\ln p_i+(1-y_i)\ln(1-p_i)]$를 최소화한다. 이 손실은 ``확률을 정직하게 말하면'' 가장 작아지므로(proper scoring rule), 로지스틱의 확률은 수준이 대체로 잘 맞는다.

\section{규제: Ridge, Lasso, Elastic Net}
변수가 100개 넘고 서로 겹치면(순위상관 0.9 이상 12쌍), 로지스틱은 학습 데이터의 우연까지 외운다. 계수 크기에 벌점을 주면 막을 수 있다.
\[
\min_\beta\ \text{log loss}+\lambda\cdot\text{벌점}(\beta),\qquad
\text{Ridge: }\|\beta\|_2^2,\quad \text{Lasso: }\|\beta\|_1,\quad \text{EN: }\alpha\|\beta\|_1+(1-\alpha)\|\beta\|_2^2 .
\]
scikit-learn은 $\lambda$ 대신 $C=1/\lambda$를 쓴다. $C$가 작을수록 규제가 세다.

\begin{figure}[H]\centering
\begin{tikzpicture}[scale=0.9,font=\small]
\begin{scope}
\draw[-{Stealth}] (-2,0)--(2.4,0) node[right]{$\beta_1$}; \draw[-{Stealth}] (0,-2)--(0,2.4) node[above]{$\beta_2$};
\draw[thick] (1.2,0)--(0,1.2)--(-1.2,0)--(0,-1.2)--cycle;
\foreach \s in {0.6,1.0,1.4} \draw[black!50,dashed,rotate around={30:(1.9,1.6)}] (1.9,1.6) ellipse ({\s*1.1} and {\s*0.55});
\fill (0,1.2) circle (2pt) node[left]{해: $\beta_1=0$};
\node at (0,-2.5) {Lasso: 모서리에 닿아 0이 생김};
\end{scope}
\begin{scope}[xshift=6.5cm]
\draw[-{Stealth}] (-2,0)--(2.4,0) node[right]{$\beta_1$}; \draw[-{Stealth}] (0,-2)--(0,2.4) node[above]{$\beta_2$};
\draw[thick] (0,0) circle (1.2);
\foreach \s in {0.6,1.0,1.4} \draw[black!50,dashed,rotate around={30:(1.9,1.6)}] (1.9,1.6) ellipse ({\s*1.1} and {\s*0.55});
\fill (0.75,0.94) circle (2pt) node[above left=1pt]{해};
\node at (0,-2.5) {Ridge: 둥글어서 0이 잘 안 생김};
\end{scope}
\end{tikzpicture}
\caption{벌점의 모양. 점선은 손실의 등고선, 굵은 선은 허용 영역. 등고선이 영역에 처음 닿는 점이 해다.}
\end{figure}

\begin{keybox}[L1과 L2]
Ridge는 모든 계수를 조금씩 줄이고, Lasso는 덜 중요한 계수를 \textbf{정확히 0}으로 만들어 변수를 고른다. Elastic Net은 둘을 섞어, 서로 겹치는 변수 묶음을 함께 남긴다.
\Intuit{Ridge는 모두의 월급을 10\% 깎고, Lasso는 몇 명을 해고한다.}
\end{keybox}

\paragraph{$C$는 어떻게 골랐나.} 평가 연도를 보면 안 되므로, 학습 데이터 안에서 20\%를 떼어 $C$ 후보 6개(0.003, 0.01, 0.03, 0.1, 0.3, 1) 중 그 20\%의 log loss가 가장 작은 값을 골랐다(\ft{choose_C}). 2013년 Lasso는 $C=0.3$, 계수 90개 중 58개만 남았다.

\section{결정나무와 앙상블}
나무는 ``FICO $<$ 690인가?'' 같은 질문을 반복해 대출을 잘게 나누고, 각 칸의 부도율을 예측으로 쓴다. 비선형, 상호작용, 결측을 자연스럽게 다루지만 나무 하나는 불안정하다. 그래서 여러 개를 합친다.
\begin{itemize}[leftmargin=*]
\item \textbf{Random Forest}: 데이터와 변수를 무작위로 뽑아 독립적인 나무 수백 개를 만들고 평균한다. 분산을 줄인다.
\item \textbf{Gradient boosting}: 나무를 순서대로 쌓되, 각 나무는 앞 나무들이 틀린 부분(손실의 기울기)을 고친다. 편향을 줄인다.
\end{itemize}
\textit{비유:} RF는 독립적인 전문가 300명의 다수결, 부스팅은 한 사람이 오답 노트를 보며 고쳐 나가는 공부다.

\begin{table}[H]\centering\small
\begin{tabularx}{\linewidth}{@{}p{28mm}L@{}}\toprule
구현 & 특징\\\midrule
LightGBM & 히스토그램 분할, 잎 단위 성장(leaf-wise). 빠르고 범주형과 결측을 직접 다룸. 이 프로젝트의 기본 모델\\
XGBoost & 깊이 단위 성장(depth-wise)이 기본, 2차 근사 손실. \ft{tree_method='hist'}로 LightGBM과 비슷한 속도\\
CatBoost & 범주형을 순서형 target statistic으로 처리, 대칭 나무. 결측을 기본적으로 ``최솟값''으로 취급\\
HistGB & scikit-learn의 LightGBM식 구현\\\bottomrule
\end{tabularx}
\caption{부스팅 구현 비교.}
\end{table}

\paragraph{LightGBM 고정 파라미터(\ft{models.LGB_PARAMS}).} 학습률 0.05, 잎 63개, 행 80\%와 열 60\% 표본추출, 잎당 최소 200건, L2 규제 5, 범주 평활 50. 나무 수는 최대 2000개이고, 학습 데이터의 15\%를 떼어 50번 연속 개선이 없으면 멈춘다(\textbf{early stopping}). 하이퍼파라미터를 평가 결과로 고르지 않은 것이 원칙이다.

\section{모델 비교 결과(23--25단계)}
조건: 시점 학습 2013--16, 트리는 시드 3개 평균, 판정 기준 사전 고정(기준 모델보다 4개 연도 중 3개 이상에서 선별 가치 0.1\%p 또는 AUC 0.002 넘게 나으면 도전자).

\begin{table}[H]\centering\small
\begin{tabular}{@{}lrrrl@{}}\toprule
모델 & AUC & 선별 가치 & Brier & 판정\\\midrule
LightGBM 기본 & 0.659 & 1.03\%p & 0.1186 & 기준, 95\% 구간 0.95--1.11\\
XGBoost & 0.655 & 1.04\%p & 0.1193 & 같음(차이 +0.01, $p=0.41$)\\
LightGBM + 피처 & 0.660 & 1.08\%p & 0.1186 & 1개 연도만 우위\\
RandomForest & \textbf{0.669} & 0.85\%p & 0.1166 & 실제로 낮음($-0.18$)\\
LightGBM 튜닝 & 0.648 & 0.86\%p & 0.1221 & 기각\\
CatBoost & 0.651 & 0.68\%p & 0.1250 & 기각(2016 붕괴)\\
Lasso + 피처 & 0.649 & 0.76\%p & 0.1203 & 선형 대표, 실제로 낮음($-0.28$)\\
Lasso & 0.652 & 0.66\%p & 0.1190 & \\
Ridge & 0.653 & 0.63\%p & 0.1189 & \\
HistGB & 0.628 & 0.63\%p & 0.1248 & 기각(2013 AUC 0.505)\\
L2 기준($C=1$) & 0.641 & 0.61\%p & 0.1198 & 선형 기준\\\bottomrule
\end{tabular}
\caption{모델 비교. ``실제로 낮음''은 25단계 짝 부트스트랩과 Holm 보정에서 유의.}
\end{table}

\begin{keybox}[이 표에서 읽을 것]
\begin{enumerate}[leftmargin=*]
\item \textbf{AUC 1위가 돈 1위가 아니다.} RF는 AUC가 가장 높지만 선별 가치는 LightGBM보다 0.18\%p 낮다(C19).
\item \textbf{1\%p는 모델의 산물이 아니다.} LightGBM과 XGBoost가 같은 값을 낸다(C20).
\item 튜닝, 피처, 다른 부스팅 구현으로 올라가지 않는다. 신호의 크기가 정해져 있다.
\end{enumerate}
\Intuit{더 똑똑한 모델을 찾는 대신, 이미 있는 신호의 크기를 정확히 쟀다.}
\end{keybox}

\subsection*{CatBoost 붕괴의 진단(24b단계)}
CatBoost는 2016년 평균 부도확률을 3.6\%로 예측했다(실제 15.9\%). 두 번 비판했다.
\begin{enumerate}[leftmargin=*]
\item \textbf{순위 문제인가 수준 문제인가?} 평균만 LightGBM 수준으로 옮기면 선별 가치가 $-0.12$에서 0.69\%p로 회복. \ensuremath{\to} 주원인은 수준.
\item \textbf{원인 가설: 새 컬럼이 ``빨리 끝난 대출''의 대리 변수였다?} 학습에서 새 컬럼 값이 있는 대출과 없는 대출의 부도율이 17.2\% vs 17.6\%로 같다. \ensuremath{\to} 가설 기각.
\end{enumerate}
가용성 규칙(\S4.5)을 적용하면 평균 PD 11.8\%, 선별 가치 0.93\%p로 정상화된다. 기제는 아직 모른다. 학습 때 소수 대출에만 있던 컬럼의 분할이 평가 때 넓게 적용되며 외삽이 틀린 것으로 보인다.

\section{기대수익 점수와 LGD}
\textbf{질문.} 부도확률 5\%짜리 8\% 금리 대출과, 25\%짜리 22\% 금리 대출 중 무엇이 나은가?

1년 동안 대출 하나에서 기대할 수 있는 초과수익을 아주 단순하게 쓰면: 확률 $1-p$로 금리 $r$를 받고, 확률 $p$로 원금의 LGD만큼 잃고, 대신 국채 $tb$를 포기한다.
\[
s=(1-p)\,r-p\cdot\mathrm{LGD}-tb .
\]
LGD는 그 시점까지 확정된 부도 대출에서 ``1 $-$ 받은 총액/대출액을 평균한 값이다. 코드는 \ft{asof.loss_given_default}이다. 실제 값은 2013년 모델 0.618, 2016년 모델 0.520이다.

\begin{workout}[두 대출의 기대수익 점수]
\Use{기대수익 점수(이 절), LGD}
\Step{1} 예시: LGD $=0.6$, 국채 $tb=1.5\%$.
\Step{2} 대출 1: $p=0.05$, $r=0.08$. $s=0.95\times0.08-0.05\times0.6-0.015=0.076-0.030-0.015=0.031$.
\Step{3} 대출 2: $p=0.25$, $r=0.22$. $s=0.75\times0.22-0.25\times0.6-0.015=0.165-0.150-0.015=0.000$.
\Ans{대출 1이 낫다(3.1\%p vs 0). 대출 2는 높은 금리가 부도 위험을 딱 상쇄하는 ``공정한 가격''이다.}
\Verify{\ft{checks.py} 5번.}
\Trap{부도확률만 보고 고른다. 대출 2의 부도확률이 30\%였다면 점수는 음수지만, 15\%였다면 대출 1보다 좋다. 금리와 함께 봐야 한다.}
\end{workout}

이 식은 연 단위 단순 근사다. 조기상환, 부도 시점, 회수 시점을 무시한다. 그래도 순위를 정하는 데는 충분하다는 것이 실험 결과다.

\section{확률 보정(calibration)}
기대수익 점수는 $p$의 \textbf{순위}뿐 아니라 \textbf{크기}에도 기댄다. $p$를 일괄로 두 배로 틀리면, 금리와 곱해지면서 등급 사이의 순서가 바뀐다.

9단계에서 보정 세 가지(isotonic, 만기 도달 대출로 보정, 생존분석 이동)를 시험했지만 모두 수준을 못 고쳤다. 과거로 만든 보정은 신용이 악화되는 추세(금리는 내리고 부도율과 DTI는 오름, C9)를 따라가지 못한다.

\begin{whybox}[그런데 왜 성과는 괜찮았나: 상쇄]
시점 학습 모델은 PD를 과소추정하고(최근 대출일수록 아직 부도가 덜 확정되어서), LGD는 과대추정한다(일찍 부도 난 대출은 받은 돈이 적어서). 두 오차가 $p\cdot\mathrm{LGD}$에서 서로 상쇄된다(C10). 2013년 LightGBM은 PD를 6.4\%로 실제(12.2\%)의 절반만 예측했지만 선별 가치는 가장 높았다. 수준은 \textbf{극단적으로} 무너질 때만(CatBoost 2016) 문제가 된다.
\end{whybox}

\begin{career}
``모델 선택을 어떻게 했나?''에 대한 좋은 답: ``판정 기준을 먼저 정하고, 같은 시점 학습 조건에서 17종을 비교했고, 이긴 모델이 없어서 원래 모델을 유지했다. 대신 왜 AUC와 돈이 어긋나는지(확률 수준), 한 모델이 왜 무너졌는지(변수 가용성)를 진단했다.'' 고른 모델보다 \textbf{고르는 절차}가 평가받는다.
\end{career}

\chapter{평가 지표: 무엇으로 채점하나}

\section{Sharpe ratio와 그 함정}
\textbf{질문.} 수익이 높지만 출렁이는 전략과, 낮지만 안정된 전략 중 무엇이 나은가?

Sharpe ratio는 초과수익 평균을 그 표준편차로 나눠 ``위험 한 단위당 보상''을 잰다: $\mathrm{SR}=\bar x/\mathrm{sd}(x)$, $x=$ 수익률 $-$ $r_f$.

과제식 Sharpe는 대출마다 $x$를 정한다: 승인하면 그 대출의 초과수익, \textbf{거절하면 국채를 산 것으로 보아 0}. 승인비율을 $a$, 승인 대출 초과수익의 평균과 표준편차를 $m,s$라 하면 전체의 평균은 $am$, 분산은 $a(s^2+m^2)-a^2m^2$이므로
\[
\mathrm{SR}(a)=\frac{am}{\sqrt{a s^2+a(1-a)m^2}}=\sqrt{a}\cdot\frac{m}{\sqrt{s^2+(1-a)m^2}} .
\]

\begin{workout}[승인비율에 따른 Sharpe]
\Use{$\sqrt a$ 항등식(이 절)}
\Step{1} 예시: 승인 대출의 $m=4\%$, $s=20\%$로 고정(같은 대출들을 고르는 비율만 바꿈).
\Step{2} $a=1$: $\mathrm{SR}=0.04/0.20=0.200$.
\Step{3} $a=0.5$: $\sqrt{0.5}\times0.04/\sqrt{0.04+0.5\times0.0016}=0.7071\times0.04/0.2020=0.140$.
\Step{4} $a=0.1$: $0.3162\times0.04/\sqrt{0.04+0.9\times0.0016}=0.0621$.
\Ans{대출의 질이 같아도 승인비율만 줄이면 Sharpe가 $\sqrt a$에 가깝게 떨어진다.}
\Verify{\ft{checks.py} 7번.}
\Trap{Sharpe가 승인비율에 따라 역U자를 그리면 ``최적 승인비율''을 찾았다고 생각한다. 그 모양의 상당 부분은 이 항등식에서 기계적으로 나온다(C4).}
\end{workout}

\begin{keybox}[비율 지표의 공통 문제(C11)]
Sharpe, information ratio, RAROC 같은 비율 지표는 분모(위험, 자본)를 줄이는 쪽, 즉 \textbf{투자를 줄이는 쪽}으로 기운다. 2013--16년처럼 전부 승인해도 국채를 이기는 시장에서는 ``얼마나 많이 벌었나''가 빠진다.
\Intuit{비율은 효율을 재지 규모를 재지 않는다. 1원 넣어 2원 번 사람이 100원 넣어 150원 번 사람보다 ``더 잘했다''고 말한다.}
\end{keybox}

\section{대안 지표들}
\begin{table}[H]\centering\small
\begin{tabularx}{\linewidth}{@{}p{30mm}LL@{}}\toprule
지표 & 식/뜻 & 이 프로젝트에서\\\midrule
Information ratio & 기준(전부 승인) 대비 초과수익 평균/표준편차 & 비율 지표의 같은 문제\\
RAROC & (이자 $-$ 조달비 $-$ 운영비 $-$ 예상손실)/바젤 자본 & 기본 가정에서 99\% 승인, 결과가 가정에 좌우. 주 지표에서 제외\\
EVA & 수익 $-$ 요구수익률 \ensuremath{\times} 자본 & RAROC과 같은 가정 의존\\
CVaR(하위 5\% 평균) & 최악 꼬리의 평균 & 보조 보고\\
확실성 등가 & 위험회피 효용으로 환산한 수익 & 위험회피 50까지도 최적은 전부 승인\\
\textbf{선별 곡선} & 승인비율별 (선택 IRR $-$ 무작위 IRR) & \textbf{주 지표}\\
\textbf{Brinson 분해} & 초과수익 = 배분 + 선별 + 상호작용 & \textbf{주 해석 도구}\\\bottomrule
\end{tabularx}
\caption{지표 비교(10, 11, 14단계).}
\end{table}

\section{순위의 정확도와 수준의 정확도: AUC와 Brier}
\textbf{AUC}는 부도 하나와 정상 하나를 무작위로 뽑았을 때 부도 쪽에 더 높은 점수를 줄 확률이다. 0.5는 동전 던지기, 1은 완벽. \textbf{Brier}는 $\frac1n\sum(p_i-y_i)^2$로, 확률 크기가 맞는지를 본다.

\begin{workout}[5건짜리 AUC와 Brier]
\Use{AUC, Brier(이 절)}
\Step{1} 예시: 부도 2건(예측 0.30, 0.10), 정상 3건(0.20, 0.05, 0.08).
\Step{2} 부도-정상 쌍 $2\times3=6$개를 센다. 0.30은 세 정상 모두보다 높다(3쌍). 0.10은 0.05, 0.08보다 높고 0.20보다 낮다(2쌍). \Concept{AUC 정의}
\Step{3} AUC $=5/6=0.833$.
\Step{4} Brier $=\frac15[(0.3-1)^2+(0.1-1)^2+0.2^2+0.05^2+0.08^2]=\frac15[0.49+0.81+0.04+0.0025+0.0064]=0.270$.
\Ans{AUC 0.833, Brier 0.270. 모든 확률을 절반으로 줄여도 AUC는 그대로지만 Brier는 바뀐다.}
\Verify{\ft{checks.py} 6번.}
\end{workout}

\section{선별 곡선과 선별 가치}
\textbf{질문.} 순위가 돈을 얼마나 더 버는지를 한 숫자로?

승인비율 $q=10\%,20\%,\dots,90\%$마다 점수 상위 $q$를 승인한 포트폴리오의 발행월 코호트 IRR에서 국채를 뺀다. 같은 비율을 무작위로 승인한 경우(5번 평균)도 계산해 뺀다. 그 차이를 9개 비율에 걸쳐 평균한 것이 \textbf{선별 가치}다.
\[
\mathrm{SV}=\frac19\sum_{q}\Big[\overline{(\mathrm{IRR}-tb)}_{\text{상위 }q}-\overline{(\mathrm{IRR}-tb)}_{\text{무작위 }q}\Big].
\]
무작위를 빼므로 ``그 해 시장이 좋았다''는 효과가 사라지고 순위의 효과만 남는다. 승인비율 하나를 고르지 않으므로 임계값 선택에 따른 엿보기도 없다.

\begin{insight}[선행 연구]
처음엔 새로운 방법이라 생각했지만, 수익 기반 점수(profit scoring, Serrano-Cinca \& Guti\'errez-Nieto 2016)와 EMP(expected maximum profit, Verbraken 외 2014)가 이미 같은 생각을 했다. 이 프로젝트의 기여는 방법이 아니라 \textbf{시점 학습 + 코호트 IRR + 부트스트랩 + 종결 편향 점검}으로 정직하게 잰 것이다.
\end{insight}

\section{Brinson 분해: 추가 수익은 어디서 오나}
펀드 평가에서 쓰는 방법이다. 기준 포트폴리오(전부 승인)와 모델 포트폴리오(상위 50\%)를 등급 $k$별로 비교한다. 기준 비중 $w^b_k$, 모델 비중 $w^p_k$, 기준 수익 $R^b_k$, 모델 수익 $R^p_k$라 하면
\[
\underbrace{\sum_k w^p_kR^p_k-\sum_k w^b_kR^b_k}_{\text{초과수익}}
=\underbrace{\sum_k(w^p_k-w^b_k)R^b_k}_{\text{배분}}
+\underbrace{\sum_k w^b_k(R^p_k-R^b_k)}_{\text{선별}}
+\underbrace{\sum_k(w^p_k-w^b_k)(R^p_k-R^b_k)}_{\text{상호작용}} .
\]

\begin{workout}[두 등급 Brinson]
\Use{Brinson 분해(이 절)}
\Step{1} 예시: 등급 A, B. 기준 비중 (0.6, 0.4), 수익 (4\%, 7\%). 모델 비중 (0.7, 0.3), 수익 (5\%, 8\%).
\Step{2} 기준 수익 $0.6\times4+0.4\times7=5.2\%$. 모델 수익 $0.7\times5+0.3\times8=5.9\%$. 초과 0.7\%p.
\Step{3} 배분 $=(0.1)(4)+(-0.1)(7)=-0.3\%$p. \Concept{배분: 등급 비중 바꾸기}
\Step{4} 선별 $=0.6(1)+0.4(1)=+1.0\%$p. \Concept{선별: 같은 등급 안에서 고르기}
\Step{5} 상호작용 $=(0.1)(1)+(-0.1)(1)=0$.
\Ans{$-0.3+1.0+0=0.7\%$p. 이 모델은 등급 배분으로는 손해를 보고, 등급 안 선별로 벌었다.}
\Verify{\ft{checks.py} 9번. 세 항의 합이 초과수익과 같다.}
\end{workout}

\begin{keybox}[결과 C7, C13, C17]
실제 데이터(세부등급 35개 기준, 시드 5개, 2017년 보류 구간)에서도 모델의 추가 수익은 주로 \textbf{등급 안 선별}에서 나왔다. 상호작용은 $-0.2$--$+0.1$\%p로 작았다. 묶음 기준을 금리 20구간, 목적, 주, 금액으로 바꾸거나 60개월 대출로 바꿔도 유지된다. 단 2013년은 배분이 주효과였다.
\Intuit{LC는 등급끼리의 가격은 잘 매겼다. 돈은 같은 등급 안의 ``숨은 차이''에서 나온다.}
\end{keybox}

\begin{career}
금융 도메인에서 Brinson attribution은 익숙한 언어다. ML 모델의 성과를 ``AUC 0.66''이 아니라 ``추가 수익 중 몇 \%p가 등급 배분, 몇 \%p가 등급 안 선별''로 설명하면, 현업(투자, 리스크) 면접관과 같은 언어로 말하게 된다.
\end{career}

\chapter{검증: 우연이 아님을 보이는 법}

``1.03\%p 더 번다''는 숫자 하나는 아무것도 증명하지 않는다. 시드가 좋았을 수도, 특정 해가 좋았을 수도, 수십 가지를 시도하다 하나가 걸렸을 수도 있다. 이 장은 그 가능성을 하나씩 지우는 도구들이다.

\section{사전 등록: 결과를 보기 전에 기준을 정한다}
11단계부터 모든 실험은 \textbf{판정 기준을 코드 맨 위 설명에 먼저 적고} 실행했다. 예: ``가장 단순한 방식보다 4개 연도 중 3개 이상에서 선별 가치 0.1\%p 또는 AUC 0.002 넘게 나으면 채택, 없으면 단순한 방식.''

\begin{keybox}[왜 중요한가]
결과를 보고 기준을 정하면, 어떤 결과든 ``성공''으로 만들 수 있다(garden of forking paths). 기준을 먼저 적어 두면 실패도 실패로 기록된다. 이 프로젝트의 실패 목록(연구 노트 2절)이 긴 것은 이 원칙의 결과다.
\Intuit{과녁을 먼저 그리고 쏴야 한다. 쏜 뒤에 화살 주위에 과녁을 그리면 백발백중이다.}
\end{keybox}

test 데이터에 대해서는 한 단계 더 나갔다. 평가 규칙, 모델, 가설, 지지 기준을 \ft{docs/test_protocol.md}에 적고 git 태그 \ft{test-protocol-v1}을 달았다. test를 본 뒤 바꿔야 하면 문서를 고치지 않고 ``이탈 기록''에 남긴다.

\section{위약 검정과 시드}
\begin{itemize}[leftmargin=*]
\item \textbf{라벨 섞기(placebo)}: 학습 라벨을 무작위로 섞어 모델을 만들면, 선별 가치는 0 근처여야 한다. 실제로 0 근처였다. 평가 코드 자체가 양수를 만들어 내는 버그가 없다는 확인이다.
\item \textbf{무작위 승인}: 같은 비율을 무작위로 승인한 포트폴리오를 기준선으로 뺀다(\S6.4).
\item \textbf{시드}: 모델 학습의 무작위성(행/열 표본추출, early stopping 분할)만 바꿔도 선별 가치가 최대 0.09\%p 움직인다. 효과가 이보다 작으면 의미가 없다(C16). 그래서 비교는 시드 3--10개 평균으로 한다.
\end{itemize}

\section{발행월 부트스트랩}
\textbf{질문.} 선별 가치 1.03\%p의 불확실성은 얼마인가?

같은 달 발행된 대출들은 같은 경기를 겪으므로 서로 독립이 아니다. 대출 단위로 다시 뽑으면 불확실성을 과소평가한다. 그래서 \textbf{발행월}을 단위로 다시 뽑는다(block bootstrap). 선별 가치는 월별 값 $f_m$의 평균이므로, 48개 월을 복원추출로 48번 뽑아 평균 내기를 1만 번 반복하고, 그 분포의 2.5\%, 97.5\% 분위를 95\% 구간으로 쓴다.

\begin{workout}[월별 차이 5개의 부트스트랩]
\Use{bootstrap(이 절)}
\Step{1} 예시: 두 모델의 월별 선별 가치 차이 $(1.2, 0.8, -0.2, 1.0, 0.6)$\%p. 평균 0.68\%p.
\Step{2} 5개 중 5개를 복원추출(예: 0.8, 0.8, 1.0, $-0.2$, 1.2)하고 평균을 낸다.
\Step{3} 1만 번 반복해 평균 1만 개를 얻는다. \Concept{bootstrap 분포}
\Step{4} 2.5\%와 97.5\% 분위: 0.24\%p와 1.04\%p.
\Ans{95\% 구간 (0.24, 1.04)\%p. 0을 포함하지 않으므로 ``차이가 있다''.}
\Verify{\ft{checks.py} 11번(시드 0).}
\Trap{두 모델을 따로 부트스트랩해 구간이 겹치는지 본다. 같은 달끼리 짝지은 \textbf{차이}를 부트스트랩해야 공통 변동이 지워져 훨씬 정밀하다(paired).}
\end{workout}

\section{여러 번 비교할 때: Holm 보정}
4개 모델을 기준과 비교하면 검정이 4번이다. 각각 5\% 수준이면 하나쯤은 우연히 ``유의''가 나올 수 있다. Holm 보정은 $p$값을 작은 순으로 정렬하고 $j$번째에 $(k-j+1)$을 곱한 뒤, 앞의 값보다 작아지지 않게 누적 최댓값을 취한다.

\begin{workout}[Holm 보정]
\Use{Holm(이 절)}
\Step{1} 예시: $k=3$, $p=(0.01, 0.02, 0.04)$.
\Step{2} $0.01\times3=0.03$, $0.02\times2=0.04$, $0.04\times1=0.04$.
\Step{3} 누적 최댓값: $(0.03, 0.04, 0.04)$.
\Ans{세 개 모두 0.05 미만이라 유의. Bonferroni(모두 \ensuremath{\times}3)였다면 0.04가 0.12가 되어 세 번째는 탈락.}
\Verify{\ft{checks.py} 12번.}
\end{workout}

\section{데이터 엿보기와 보류 구간}
2013--16년 평가 구간을 7단계부터 26단계까지 계속 다시 썼다. 규칙, 지표, 피처, 모델을 이 구간에서 고를수록 이 구간에 맞춰질 위험이 커진다. 지금까지 비교한 규칙, 지표, 모델 조합은 약 80개다.

\begin{itemize}[leftmargin=*]
\item 12단계 EDA는 2007--12년만 봤다.
\item 14단계에서 처음으로 2017년 발행분을 보류 구간으로 썼다(선별 가치 1.22\%p, 다만 종결 편향 있음). 이후 이 구간도 ``본 구간''이다.
\item 23--26단계 결과가 좋아도 test 프로토콜의 모델은 바꾸지 않았다.
\item 최종 판단은 한 번도 보지 않은 공식 test로만 한다.
\end{itemize}

\begin{warnbox}[test 해석 규칙(26단계에서 미리 정함)]
test에도 진행 중 대출이 없으므로 종결 편향이 생긴다. 발행연도별 만기 도달 비율부터 확인하고, 비율이 높은 연도만 기대수익 순위 vs 부도확률 순위 비교의 근거로 쓴다. 비율이 낮은 연도에서 부도확률 순위가 이기면 ``예상된 편향 방향''으로 보고하고, 그것만으로 C6를 기각하지 않는다.
\end{warnbox}

\section{주장 장부: 확신의 등급}
모든 주장은 \ft{docs/research_log.md}의 장부에 근거 파일, 상태(지지/잠정/기각), 남은 약점과 함께 적는다. 예:

\begin{table}[H]\centering\small
\begin{tabularx}{\linewidth}{@{}p{8mm}Lp{20mm}L@{}}\toprule
\# & 주장 & 상태 & 남은 약점\\\midrule
C6 & 기대수익 순위가 부도확률 순위보다 더 번다 & 지지(전체 표본) & 격차가 해마다 줄어듦\\
C6b & 부도확률 순위는 무작위와 비슷하다 & \textbf{기각} & 시드 5개로 보면 유의하게 양수\\
C7 & 가치는 등급 안 선별에서 나온다 & 지지 & 2013년은 배분이 주효과\\
C19 & AUC 순위는 돈 순위와 다르다 & 지지(탐색) & 수준은 극단적일 때만 문제\\
C23 & 종결 편향은 부도확률 순위를 유리하게 만든다 & 지지 & 2013--15년 36개월만\\\bottomrule
\end{tabularx}
\caption{주장 장부 발췌.}
\end{table}

\begin{career}
``결과가 우연이 아니라는 걸 어떻게 아나요?''에 대한 답이 이 장 전체다. 사전 등록, 위약 검정, 다중 시드, 군집 부트스트랩, 다중비교 보정, 보류 구간, 그리고 ``틀린 주장을 기각한 기록''. 특히 마지막이 신뢰를 만든다.
\end{career}

\chapter{코드 구조와 작업 방식}

\section{저장소 지도}
\begin{table}[H]\centering\small
\begin{tabularx}{\linewidth}{@{}p{38mm}L@{}}\toprule
경로 & 역할\\\midrule
\ft{src/lc/config.py} & 경로 상수(데이터, 국채 CSV, 결과 폴더)\\
\ft{src/lc/data.py} & 원본 읽기, 누수 컬럼 삭제, 라벨(\ft{bad}, \ft{resolved}) 만들기, 특성 101개\\
\ft{src/lc/returns.py} & 현금흐름 복원, IRR, 네 가지 수익률, 국채 붙이기\\
\ft{src/lc/asof.py} & 시점 학습 표본(\ft{known_at}, \ft{training_set}), LGD\\
\ft{src/lc/models.py} & LightGBM 고정 파라미터와 학습 함수\\
\ft{src/lc/evaluate.py} & 코호트 IRR, Sharpe, 요약 지표\\
\ft{src/lc/missing.py} & 결측 유형 분류와 처리 방식 5종(\ft{MissingPolicy})\\
\ft{src/lc/woe.py} & WOE 변환과 IV 표\\
\ft{src/lc/features.py} & 엔지니어링 피처, out-of-fold 우편번호 인코딩\\
\ft{src/lc/calibration.py}, \ft{raroc.py}, \ft{benchmarks.py} & 보정, RAROC, 비교 기준\\
\ft{scripts/02_...}--\ft{26_...} & 단계별 실험. 번호 = 커밋 단계\\
\ft{tests/} & 손계산 테스트 27개(\ft{pytest tests})\\
\ft{results/} & 각 실험의 CSV와 로그\\
\ft{docs/} & 결정 기록, 연구 노트, test 프로토콜, 이 노트\\\bottomrule
\end{tabularx}
\caption{저장소 구조.}
\end{table}

\section{핵심 코드 읽기}
\subsection*{시점 학습 표본}
\begin{lstlisting}
def known_at(d, cutoff, lag=CHARGE_OFF_LAG):
    """Loans whose outcome was settled by cutoff."""
    age = months_to(d, cutoff)                 # months from issue to cutoff
    settle = d['tau'].to_numpy() + np.where(d['bad'].to_numpy() == 1, lag, 0)
    return (age > 0) & d['resolved'].to_numpy() & (settle <= age)
\end{lstlisting}
한 줄씩: \ft{age}는 기준일에 대출이 몇 개월 되었나. \ft{settle}은 결과가 확정되는 개월(정상은 마지막 납부, 부도는 +5). 기준일 전에 발행되었고(\ft{age > 0}), 지금은 끝났고(\ft{resolved}), 그 끝이 기준일 전에 확정되었으면(\ft{settle <= age}) 학습에 쓴다. 이 세 줄이 3장 그림의 A, B를 남기고 C, D를 빼는 일을 한다.

\subsection*{코호트 IRR}
\begin{lstlisting}
def cohort_irr(d, approve, fee=0.01):
    a = np.asarray(approve, bool)
    rows = []
    for ym, idx in d.groupby('ym').indices.items():   # each issue month
        sel = idx[a[idx]]                              # approved loans of that month
        tb = d['tb'].to_numpy()[idx].mean()            # treasury of that month
        v = returns.pooled_irr(d.iloc[sel], fee=fee) if len(sel) else tb
        rows.append((ym, v, tb, len(sel) / len(idx)))
    return pd.DataFrame(rows, columns=['ym', 'irr', 'tb', 'approve_rate']).set_index('ym')
\end{lstlisting}
그 달에 하나도 승인하지 않으면 국채를 산 것으로 본다(\ft{v = tb}). 그래서 거절은 초과수익 0이다(\S6.1).

\subsection*{기대수익 점수와 선별 곡선}
\begin{lstlisting}
p = model.predict_proba(X_eval)[:, 1]                      # default probability
score = (1 - p) * ev['r'] - p * lgd - ev['tb']            # expected-return score
QS = [0.1, 0.2, ..., 0.9]
curve = [ (cohort_irr(ev, score >= np.quantile(score, 1 - q))
           .eval('irr - tb').mean()) for q in QS ]          # approve top q
sv = np.mean(np.array(curve) - random_curve)               # selection value
\end{lstlisting}

\subsection*{선형 모델 전처리(가용성 규칙)}
\begin{lstlisting}
avail = X_train.notna().mean()
cols = avail[avail >= 0.5].index                 # 1) keep features >50% present in training
lo, hi = X_train[cols].quantile(0.01), X_train[cols].quantile(0.99)
X = X[cols].clip(lo, hi, axis=1)                  # 2) clip at training 1%/99%
X = X.fillna(X_train[cols].median())              # 3) fill with training median
\end{lstlisting}
모든 기준값(분위, 중앙값)은 \textbf{학습 데이터에서만} 계산한다. 평가 데이터로 계산하면 그것도 누수다.

\section{테스트}
\ft{tests/}의 27개 테스트는 손으로 계산할 수 있는 작은 예로 함수를 확인한다. 예: 정상 상환 대출의 IRR이 약정 금리(월 복리)와 같은지, \ft{known_at}이 기준일 뒤 확정된 부도를 빼는지, 결측 처리 값이 학습 데이터에서만 오는지. 숫자 실험을 바꿀 때마다 \ft{pytest tests}를 돌린다.

\section{작업 방식: 단계, 커밋, 기록}
\begin{enumerate}[leftmargin=*]
\item 질문과 판정 기준을 스크립트 맨 위에 쓴다.
\item 실행하고 결과를 \ft{results/}에 저장한다.
\item 비판한다(우연인가, 누수인가, 다른 설명은?). 필요하면 진단 실험을 한다(예: 24b).
\item \ft{docs/decisions.md}(무엇을 왜 골랐나)와 \ft{docs/research_log.md}(주장, 실패, 정정, 가정, 엿보기)를 갱신한다.
\item 한 단계 = 한 커밋. 실험은 브랜치에서 하고 main에는 바로 올리지 않는다.
\end{enumerate}

\begin{career}
이 구조는 연구 재현성(reproducibility)의 기본형이다: 모듈화된 라이브러리 + 번호 붙은 실험 스크립트 + 결과 파일 + 결정 로그 + 테스트. ``노트북 하나에 다 했다''와 비교하면, 누가 어떤 숫자를 물어도 어느 커밋의 어느 스크립트에서 나왔는지 답할 수 있다.
\end{career}

\chapter{전체 이야기: 팀과 교수님께 설명하기}

\section{10분 발표의 뼈대}
\begin{enumerate}[leftmargin=*]
\item \textbf{문제 재정의(1분).} ``Sharpe를 최대화하라''는 과제를 ``투자자가 무엇을 사야 하고, 그것이 운이 아님을 어떻게 보이나''로 다시 쓴다. 이유: Sharpe는 거절을 0으로 두면 승인비율에 따라 기계적으로 변한다($\sqrt a$ 항등식).
\item \textbf{자를 먼저 정한다(2분).} 수익률 정의 네 가지가 결론을 뒤집는다. 현금흐름에서 복원한 코호트 IRR(수수료 1\%)과 만기를 맞춘 국채를 쓴다. 이 자로 보면 전부 승인도 국채를 이긴다.
\item \textbf{시간을 지킨다(2분).} 시점 학습(상각 지연 5개월 포함). 무작위 분할의 AUC 0.73은 부풀려진 값이고 실제는 0.62--0.68.
\item \textbf{핵심 결과(2분).} 기대수익 순위는 무작위보다 1.03\%p(0.95--1.11) 더 번다. 추가 수익은 같은 등급 안의 선별에서 나온다. LC는 신용점수는 잘 가격에 넣었지만 상환 능력(소득, DTI)은 덜 넣었다.
\item \textbf{견고성(2분).} 모델 17종, 결측 처리 9종, 피처 6종을 사전 기준으로 비교: 아무것도 1\%p를 바꾸지 못했다. XGBoost도 같은 값. AUC 1위(RF)가 돈 1위가 아니다.
\item \textbf{정직성(1분).} 틀렸던 주장과 기각한 가설 목록, 종결 편향 시뮬레이션, test 프로토콜 사전 고정.
\end{enumerate}

\section{예상 질문과 답}
\begin{qa}{왜 Sharpe를 안 쓰나?}
쓰지만 주 지표로 두지 않는다. 거절을 무위험 0으로 두면 Sharpe가 승인비율의 제곱근에 끌려간다. 또 비율 지표는 규모를 무시해 투자를 줄이는 쪽으로 기운다. 그래서 승인비율 전체에 걸친 선별 곡선을 주 지표로, Sharpe는 과제 요구에 맞춰 함께 보고한다.
\end{qa}
\begin{qa}{AUC가 0.66밖에 안 되는데 의미가 있나?}
과제는 맞히기가 아니라 돈이다. AUC 0.66으로도 무작위보다 연 1\%p를 더 벌고, 그 구간이 0을 멀리 벗어난다. AUC가 더 높은 RF는 오히려 덜 번다.
\end{qa}
\begin{qa}{2013--16년을 너무 많이 본 것 아닌가?}
맞다. 그래서 (1) 판정 기준을 미리 정했고, (2) 결과가 좋아도 test 모델을 바꾸지 않았고, (3) 시도한 조합 수(약 80개)를 공개하고, (4) 최종 판단은 test로만 한다.
\end{qa}
\begin{qa}{test에서 효과가 사라지면?}
그것도 결론이다. 소득 효과가 해마다 줄었고(2.66 \ensuremath{\to} 0.19\%p), 선별 가치도 줄었다. 사라지면 ``LC 가격이 점점 효율적이 되었다''가 결론이 된다. 어느 쪽이든 미리 정한 기준으로 보고한다.
\end{qa}
\begin{qa}{우승 논문과 무엇이 다른가?}
그 논문은 단순 연환산과 무작위 분할에 가까운 설정을 썼다. 우리는 같은 데이터로 그 설정이 ``전부 승인 < 국채''라는 결론을 만들어 내는 경로를 재현했고(C3, 잠정), 현금흐름 IRR과 시점 학습으로 다시 쟀다.
\end{qa}

\section{남은 한계}
\begin{itemize}[leftmargin=*]
\item 현금흐름 시점은 근사다(월별 납입 기록 없음).
\item 기대수익 점수는 조기상환과 부도 시점을 무시한 1년 근사다.
\item 등급 안 정보의 크기가 해마다 줄었다. 미래에도 있을지는 모른다.
\item 소득은 자기 신고 값이다.
\item CatBoost 붕괴의 정확한 기제는 확인하지 못했다.
\item 선별 곡선의 아이디어 자체는 선행 연구가 있다.
\end{itemize}

\section{다음 단계}
\begin{enumerate}[leftmargin=*]
\item 공식 test를 받으면 \ft{05_inspect_test.py}로 구성(발행연도, 만기 도달 비율)부터 확인한다.
\item \ft{15_test_evaluate.py}를 프로토콜 그대로 한 번 돌린다. 가용성 규칙 LightGBM은 민감도로 함께 보고한다.
\item H1--H4를 미리 정한 기준으로 ``지지/지지되지 않음''으로 보고한다.
\end{enumerate}

\appendix
\chapter{면접 Q\&A}

\begin{qa}{이 프로젝트를 2분 안에 설명해 보세요.}
LendingClub 대출 175만 건으로 투자자의 승인 전략을 만들었다. 정확도 대신 현금흐름에서 복원한 IRR을 목표로 삼았고, 결정 시점에 알 수 있던 정보로만 학습하는 시점 학습을 썼다. 기대수익 순위가 무작위보다 연 1.03\%p(95\% 구간 0.95--1.11) 더 벌었고, Brinson 분해로 그 수익이 같은 등급 안의 선별에서 나온다는 것을 보였다. 모델 17종을 사전 기준으로 비교했지만 이 값을 바꾸지 못했고, 평가 표본의 종결 편향을 시뮬레이션으로 정량화했다.
\end{qa}

\begin{qa}{Data leakage를 어떻게 막았나요?}
컬럼 수준: 승인 시점에 관측 가능한지로 판별해 상환 결과와 사후 갱신 컬럼을 뺐다. 행 수준: 연도 $Y$ 모델은 $Y$년 1월 1일까지 결과가 확정된 대출로만 학습했고, 부도 확정에 상각 지연 5개월을 넣었다. 전처리 수준: 결측 대치 값, 분위 자르기, target encoding을 모두 학습 데이터에서만(target encoding은 out-of-fold로) 계산했다.
\end{qa}

\begin{qa}{왜 무작위 분할을 쓰지 않았나요?}
대출 결과는 경기와 LC 심사 기준에 따라 시간이 흐르며 변한다. 무작위 분할은 미래 대출로 과거를 맞히게 해 AUC를 0.73까지 부풀렸고, 시점 학습에서는 0.62--0.68이었다.
\end{qa}

\begin{qa}{AUC와 실제 수익이 왜 다른가요?}
AUC는 순위만 본다. 기대수익 $(1-p)r-p\,\mathrm{LGD}$는 $p$의 크기가 금리와 곱해지므로 등급 간 순서가 확률 수준에 달려 있다. Random Forest는 AUC 1위였지만 선별 가치는 LightGBM보다 0.18\%p 낮았다(부트스트랩으로 유의).
\end{qa}

\begin{qa}{Ridge와 Lasso의 차이는?}
Ridge는 L2 벌점으로 계수를 고르게 줄이고, Lasso는 L1 벌점으로 일부 계수를 정확히 0으로 만들어 변수 선택을 한다. L1 허용 영역은 모서리가 있는 마름모라 손실 등고선이 축 위의 모서리에 먼저 닿기 때문이다. 이 프로젝트에서는 $C$를 학습 데이터 안 검증 log loss로 골랐고, 둘의 차이는 잡음 안이었다.
\end{qa}

\begin{qa}{결측을 어떻게 처리했나요?}
먼저 결측의 이유를 분류했다(구조적, 사건 없음, 시기에 따른 MAR, MNAR 의심). 트리는 결측을 그대로 두는 것이 가장 나았다. 선형 모델은 결측 처리보다 ``학습 때 거의 없던 변수''가 문제여서, 가용성 규칙(50\% 미만 제외) + 분위 자르기 + 중앙값으로 정했다. 이 규칙 없이 2016년 AUC가 0.508까지 무너졌다.
\end{qa}

\begin{qa}{결과가 우연이 아님을 어떻게 보였나요?}
사전 등록된 판정 기준, 라벨 섞기 위약 검정, 다중 시드, 발행월 단위 짝 부트스트랩, Holm 다중비교 보정, 보류 구간, 그리고 test 프로토콜의 git 태그 고정.
\end{qa}

\begin{qa}{가장 크게 틀렸던 판단은?}
두 가지. 첫째, 시드 1개 실험으로 ``부도확률 순위는 무작위와 비슷하다''고 했다가 다중 시드와 부트스트랩으로 기각했다. 둘째, 2017년 보류 구간에서 부도확률 순위가 이긴 것을 약점으로 봤는데, 이미 끝난 대출을 일부러 잘라 보니 같은 역전이 재현되어 종결 편향의 산물임을 알았다.
\end{qa}

\begin{qa}{LightGBM의 early stopping은 어떻게 했나요?}
학습 데이터의 15\%를 무작위로 떼어 검증 log loss가 50번 연속 개선되지 않으면 멈췄다. 평가 연도는 쓰지 않았다. 떼는 표본이 시드에 따라 달라 시드 간 변동의 일부가 된다.
\end{qa}

\begin{qa}{현업이라면 무엇을 더 하겠습니까?}
조기상환 위험을 별도 모델로(경쟁 위험 생존모형), LGD를 대출 특성으로 예측, 거시 변수(실업률)로 신용 사이클 반영, 그리고 실시간 운영 시 변수 가용성 모니터링(새 변수가 생기면 자동으로 가용성 규칙 적용).
\end{qa}

\chapter{Python 실습 코드}
\section{예제 검산}
본문 예제의 모든 숫자는 \ft{docs/notes/lecture/code/checks.py}로 다시 계산할 수 있다.
\begin{lstlisting}
import numpy as np
from scipy.optimize import brentq

def irr(cf):
    f = lambda m: (cf / (1 + m) ** np.arange(len(cf))).sum()
    m = brentq(f, -0.95, 1.0)           # monthly rate that zeroes the NPV
    return (1 + m) ** 12 - 1            # annualize with monthly compounding

P, r, T = 10000, 0.12, 36
i = r / 12
inst = P * i / (1 - (1 + i) ** -T)      # 332.14
print(irr(np.r_[-P, np.full(T, inst)])) # 0.1268

cf = np.zeros(16); cf[0] = -P; cf[1:11] = inst; cf[15] = 500
print(irr(cf))                          # -0.768 (default after 10 payments)
\end{lstlisting}

\section{실험 재현 순서}
\begin{lstlisting}[language=bash]
python scripts/prepare_data.py          # raw CSV -> parquet
pytest tests                            # 27 hand-checked tests
python scripts/03_returns.py            # four return definitions
python scripts/14_validate_methods.py   # selection value, Brinson, bootstrap
python scripts/23_models_final.py       # model comparison (about 9 min)
python scripts/25_paired_bootstrap.py   # paired monthly bootstrap
python scripts/26_resolution_bias.py    # resolution-bias simulation
\end{lstlisting}

\end{document}
